\documentclass[a4paper, 11pt]{report}

% ─── Packages ─────────────────────────────────────────────────────────────────
\usepackage[T1]{fontenc}
\usepackage{lmodern}
\usepackage[margin=1in]{geometry}
\hfuzz=13pt  % suppress overfull hbox warnings under 13pt (nested lists)
\usepackage{amsmath, amssymb, amsthm}
\usepackage{graphicx}
\usepackage{hyperref}
\usepackage{enumitem}
\usepackage{booktabs}
\usepackage{multirow}
\usepackage{listings}
\usepackage{xcolor}
\usepackage{fancyhdr}
\usepackage{tocloft}

% ─── Hyperref setup ───────────────────────────────────────────────────────────
\hypersetup{
    colorlinks=true,
    linkcolor=blue!70!black,
    urlcolor=blue!60!black,
    citecolor=green!50!black,
    pdftitle={Quant101 — A Quantitative Trader's Learning Lab},
    pdfauthor={Jerry Hong}
}

% ─── Fix fancyhdr warning ────────────────────────────────────────────────────
\setlength{\headheight}{14pt}
\addtolength{\topmargin}{-2pt}

% ─── Code listing style ──────────────────────────────────────────────────────
\definecolor{codebg}{RGB}{248,248,248}
\definecolor{codeframe}{RGB}{200,200,200}
\lstset{
    language=Python,
    basicstyle=\ttfamily\small,
    backgroundcolor=\color{codebg},
    frame=single,
    rulecolor=\color{codeframe},
    breaklines=true,
    showstringspaces=false,
    keywordstyle=\color{blue!70!black}\bfseries,
    commentstyle=\color{green!50!black}\itshape,
    stringstyle=\color{red!60!black},
    numbers=left,
    numberstyle=\tiny\color{gray},
    tabsize=4,
}

% ─── Theorem environments ────────────────────────────────────────────────────
\theoremstyle{definition}
\newtheorem{definition}{Definition}[chapter]
\newtheorem{example}{Example}[chapter]
\newtheorem{labmod}{Lab Module}[chapter]
\theoremstyle{remark}
\newtheorem{remark}{Remark}[chapter]
\newtheorem{experiment}{Experiment}[chapter]
\newtheorem{pitfall}{Common Pitfall}[chapter]

% ─── Header/Footer ───────────────────────────────────────────────────────────
\pagestyle{fancy}
\fancyhf{}
\fancyhead[L]{\leftmark}
\fancyhead[R]{Quant101}
\fancyfoot[C]{\thepage}

% ─── Title ────────────────────────────────────────────────────────────────────
\title{
    \textbf{Quant101} \\[0.5em]
    \Large A Quantitative Trader's Learning Lab \\[1em]
    \normalsize From Beginner to Job-Market-Ready Quant Researcher
}
\author{Jerry Hong \\ {\small with AI-assisted authoring}}
\date{Started September 2025 \\ Last updated: \today}

\begin{document}
\maketitle

% ─── Preface ──────────────────────────────────────────────────────────────────
\chapter*{Preface}
\addcontentsline{toc}{chapter}{Preface}

This document is a \textbf{learning lab notebook} --- not a textbook, not API
documentation.  It records my path from a beginning trader to a quant researcher,
maintained together with an AI agent who serves as a pair-programming partner and
knowledge editor.

The companion codebase, \texttt{quant101}, started in September~2025 as a
backtesting experiment built on Polygon.io\footnote{Rebranded as Massive around
October~2025: \url{https://massive.com/blog/polygon-is-now-massive}} flat-file
data.  It has since grown into a data pipeline, backtest engine, strategy
framework, and live-monitor system.

\subsection*{What This Document Is}

\begin{itemize}
    \item A \textbf{guided build sequence}: each chapter asks you to build
          something, then measure it, then learn the theory that explains the
          measurement.
    \item A \textbf{depth-first} exploration: fewer topics, deeper treatment,
          real experiments on real data.
    \item A \textbf{job-market bridge}: the skill progression follows what quant
          interviews actually test --- not academic curriculum order.
\end{itemize}

\subsection*{What This Document Is Not}

\begin{itemize}
    \item Not a reference manual for current file paths or API signatures --- those
          change with every refactor.
    \item Not a textbook that front-loads 200 pages of theory before any code.
    \item Not a survey of every quant topic --- Wikipedia exists.
\end{itemize}

\subsection*{How to Read}

Every chapter follows the same \textbf{lab module} structure:

\begin{enumerate}
    \item \textbf{Motivation} --- Why this matters for the project and for interviews.
    \item \textbf{Concepts} --- Just enough theory, introduced when you need it.
    \item \textbf{Implementation} --- What to build, described at the interface level
          (not hardcoded paths --- the architecture evolves).
    \item \textbf{Experiment} --- Run it on real data.  Record what you find.
    \item \textbf{Reflection} --- What worked, what broke, what to read next.
\end{enumerate}

Chapters are ordered by \textbf{build dependency}: you can read straight through
and each chapter assumes only the previous ones.  Mathematical foundations live
in Appendix~\ref{app:math} and are cross-referenced when needed --- flip there
when a chapter says ``see Appendix~A.''

\subsection*{Skill Progression}

\begin{center}
\begin{tabular}{@{}lll@{}}
\toprule
\textbf{Level} & \textbf{Parts} & \textbf{You Can\ldots} \\
\midrule
Beginner        & I, II       & Build a data pipeline, run a backtest, measure performance \\
Intermediate    & III         & Construct \& evaluate alpha factors, think cross-sectionally \\
Advanced        & IV, V       & Build portfolios, model risk, apply ML with proper validation \\
\bottomrule
\end{tabular}
\end{center}

Part~III is the \textbf{inflection point}: where a ``trader who codes'' becomes a
``quant researcher.''  The document spends the most depth there.

\tableofcontents


% ══════════════════════════════════════════════════════════════════════════════
% PART I: DATA ENGINEERING
% ══════════════════════════════════════════════════════════════════════════════

\part{Data Engineering}
\label{part:data}

\begin{quote}
\itshape
All research starts from clean data.  No analysis is trustworthy if the data
pipeline is broken.  This part covers the methodology of acquiring, cleaning, and
organizing financial market data --- the foundation everything else rests on.
\end{quote}


% ──────────────────────────────────────────────────────────────────────────────
\chapter{Data Pipeline Design}
\label{ch:pipeline}

% ── Motivation ────────────────────────────────────────────────────────────────
\section{Motivation}

Most quant tutorials skip straight to ``load a CSV and compute moving averages.''
In practice, 60--80\% of a quant researcher's time is data work: downloading,
parsing, validating, adjusting, and caching.  If you can't trust your data, every
signal you build and every backtest you run is unreliable.

This chapter covers the \emph{design} of a market data pipeline --- the stages,
the tradeoffs, and the failure modes --- without prescribing a fixed directory
layout (that evolves with the codebase).

\textbf{Interview relevance}: ``Walk me through how you'd build a data pipeline
for a factor research platform'' is a common quant infrastructure question.

% ── Concepts ──────────────────────────────────────────────────────────────────
\section{Concepts}

\subsection{The Four-Stage Pipeline}

Every market data system, whether at a hedge fund or on a personal laptop, follows
the same logical stages:

\begin{enumerate}
    \item \textbf{Acquisition} --- Download raw data from external sources.
    \item \textbf{Normalization} --- Parse, schema-map, and convert to an efficient
          columnar format.
    \item \textbf{Enrichment} --- Apply corporate actions (splits, dividends), map
          identifiers, filter the universe.
    \item \textbf{Serving} --- Load clean, adjusted data into research or backtest
          code with minimal latency.
\end{enumerate}

\begin{remark}
    The four stages are a logical decomposition.  In a small project they may
    live in two scripts; at a fund they may be four separate services with
    message queues between them.  The \emph{boundaries} matter more than the
    implementation.
\end{remark}

\subsection{Columnar Storage}

\begin{definition}[Columnar Format]
    A storage layout where values for each column are stored contiguously on disk.
    Examples: Apache Parquet, Apache Arrow, ORC\@.

    For OHLCV data with $N$ rows and 6 columns, a row-oriented CSV reads all 6
    columns even when you only need \texttt{close}.  Parquet reads only the
    \texttt{close} column --- typically 5--10$\times$ faster for analytical queries.
\end{definition}

\begin{remark}
    This project uses Parquet as the canonical on-disk format and Polars
    (which reads Parquet natively via Arrow) as the in-memory engine.
    The choice is deliberate: Polars' lazy evaluation and predicate pushdown
    make it far more efficient than Pandas for the scan-filter-aggregate pattern
    that dominates quant data work.
\end{remark}

\subsection{Incremental vs.\ Full Refresh}

Two update strategies:

\begin{itemize}
    \item \textbf{Full refresh}: Re-download and re-process everything.  Simple
          but wasteful.  Acceptable for small datasets or monthly rebuilds.
    \item \textbf{Incremental update}: Only fetch new data since the last run.
          Requires tracking watermarks (last processed date).  More complex
          but essential when daily data grows to millions of rows.
\end{itemize}

The project uses incremental updates for daily aggregates (fetch only recent
$N$ days) and full refresh for slowly-changing reference data (ticker lists,
split events).

% ── Implementation ────────────────────────────────────────────────────────────
\section{Implementation Sketch}

The pipeline is implemented as a sequence of idempotent scripts, orchestrated by
a shell-based task runner.  The key design decisions:

\begin{itemize}
    \item \textbf{Configuration-driven}: A single YAML config file
          (\texttt{basic\_config.yaml}) defines the data root directory, update
          mode (standalone / server / client), and sync settings.  The config is
          gitignored; a committed \texttt{.example} template documents the schema.
    \item \textbf{Idempotent scripts}: Each script can be re-run safely.
          Re-downloading an already-downloaded date is a no-op; re-converting an
          already-converted file overwrites with identical output.
    \item \textbf{Mode-aware orchestration}: A single entry-point script
          detects the update mode and runs the appropriate task subset
          (e.g., a client skips API fetches and syncs metadata from a server
          instead).
\end{itemize}

\begin{labmod}
    \textbf{Build it}: Set up the pipeline end-to-end for one asset class
    (e.g., US equities daily aggregates).  Download $\to$ convert $\to$ verify
    row counts $\to$ load into a Polars DataFrame and compute a 20-day SMA\@.
    Time each stage.  Where is the bottleneck?
\end{labmod}

% ── Experiment ────────────────────────────────────────────────────────────────
\section{Experiment: Pipeline Profiling}

\begin{experiment}
    \textbf{Goal}: Profile each stage of the pipeline for 7 days of US equity
    daily aggregates ($\sim$10{,}000 tickers $\times$
    7~days).

    \textbf{Method}: Time each stage separately --- download, CSV.gz
    $\to$~Parquet conversion, and Polars lazy scan + collect.

    \textbf{Record}: Wall-clock time per stage, file sizes (CSV.gz vs.\ Parquet),
    memory usage during conversion.

    \textbf{Reflection space}: \emph{(fill in after running)}
\end{experiment}


% ──────────────────────────────────────────────────────────────────────────────
\chapter{Data Quality and Corporate Actions}
\label{ch:data-quality}

% ── Motivation ────────────────────────────────────────────────────────────────
\section{Motivation}

Raw market data is \emph{never} clean.  Stock splits change historical prices,
tickers get reused by different companies, and data vendors introduce errors.
Ignoring these issues doesn't make them go away --- it makes your backtest return
a fantasy number.

This is the unglamorous foundation that separates a toy project from a usable
research platform.

\textbf{Interview relevance}: ``How do you handle stock splits in your data?''
and ``What is survivorship bias at the data level?'' are standard screening
questions.

% ── Concepts ──────────────────────────────────────────────────────────────────
\section{Concepts}

\subsection{Stock Splits and Adjustment}

\begin{definition}[Split Adjustment]
    A stock split with ratio $r$ (e.g., $r = 2$ for a 2-for-1 split) requires
    adjusting all historical prices before the split date:
    \[
        P_{\text{adj}} = P_{\text{raw}} \cdot \frac{1}{r}
    \]
    Volume is adjusted inversely: $V_{\text{adj}} = V_{\text{raw}} \cdot r$.

    Without adjustment, a 2-for-1 split looks like a 50\% crash in the price
    series, which poisons every downstream calculation.
\end{definition}

\begin{pitfall}
    \textbf{Split data has errors.}  Vendor-provided split tables contain
    incorrect ratios, missing events, and duplicate entries.  The project
    addresses this with a CSV-based error correction layer: a manually curated
    file lists split IDs to remove (erroneous) and add (missing), applied as a
    filter before adjustment.  This pattern --- \emph{override layer on top of
    vendor data} --- is standard practice at production quant firms.
\end{pitfall}

\subsection{Ticker Identity Over Time}

\begin{definition}[FIGI (Financial Instrument Global Identifier)]
    A 12-character, alphanumeric identifier assigned to financial instruments.
    Unlike ticker symbols, FIGIs are globally unique and persistent across name
    changes, mergers, and exchange transfers.
\end{definition}

The ticker identity problem: company~A trades as \texttt{AAPL}, gets acquired,
and later company~B is listed under the same symbol \texttt{AAPL}.  Using ticker
strings alone conflates two unrelated entities.

The solution is a \textbf{connected-component algorithm} on FIGI fields
(\texttt{composite\_figi} and \texttt{share\_class\_figi}).  Two ticker records
sharing any FIGI are the same instrument.  The algorithm runs on the full ticker
universe and produces a canonical mapping from every historical ticker string to
a unique entity ID\@.

\subsection{The Low-Volume Ticker Problem}

\begin{experiment}
    \textbf{Observation}: Many tickers show extended zero-volume periods
    ($>50$ consecutive trading days).  Investigation revealed these are mostly
    relisted instruments or ticker-symbol reuse by shell companies.

    \textbf{Decision}: Skip these tickers in backtest universes.  This is a
    documented data quality tradeoff --- not a silent exclusion.

    \textbf{Lesson}: Every universe filter introduces a bias.  Documenting
    \emph{why} you exclude something is as important as excluding it.
\end{experiment}

% ── Implementation ────────────────────────────────────────────────────────────
\section{Implementation Sketch}

\begin{itemize}
    \item \textbf{Split adjustment}: Applied during the data-serving stage, not
          during conversion.  This preserves raw data fidelity and allows
          re-adjustment when corrections are added.
    \item \textbf{Error correction}: A pair of CSV files per data type
          (e.g., \texttt{splits\_error.csv}) with columns
          \texttt{[id, error\_type, ...]}.  The loader reads vendor data, filters
          by the error file, then serves the corrected result.
    \item \textbf{Ticker mapping}: The FIGI connected-component algorithm runs
          once per universe refresh and caches the mapping.  Downstream code
          references entity IDs, not raw ticker strings.
\end{itemize}

\begin{labmod}
    \textbf{Build it}: Take a known split event (e.g., NVDA 10-for-1, June~2024).
    Load raw and adjusted price series.  Plot both on the same chart.  Verify the
    adjustment formula by hand for one data point.
\end{labmod}

% ── Experiment ────────────────────────────────────────────────────────────────
\section{Experiment: Split Adjustment Validation}

\begin{experiment}
    \textbf{Goal}: Validate the split adjustment pipeline against a known event.

    \textbf{Method}: For NVDA's 10-for-1 split (2024-06-10):
    \begin{enumerate}
        \item Load unadjusted close prices around the split date.
        \item Apply the split adjustment.
        \item Compare adjusted close with the value from a trusted source
              (e.g., Yahoo Finance adjusted close).
        \item Compute the percentage difference --- it should be $< 0.01\%$.
    \end{enumerate}

    \textbf{Reflection space}: \emph{(fill in after running)}
\end{experiment}


% ──────────────────────────────────────────────────────────────────────────────
\chapter{Universe Construction}
\label{ch:universe}

% ── Motivation ────────────────────────────────────────────────────────────────
\section{Motivation}

``What stocks do you include in your backtest?'' sounds simple.  It isn't.  The
choice of universe determines whether your results generalize or are an artifact
of selection.

A universe that is too narrow (only FAANG stocks) overfits to a handful of names.
A universe that is too broad (all 15{,}000 US tickers including OTC penny stocks)
drowns signals in noise and introduces illiquidity artifacts.

\textbf{Interview relevance}: ``How did you construct your backtest universe?'' is
a question that separates candidates who understand survivorship bias from those
who don't.

% ── Concepts ──────────────────────────────────────────────────────────────────
\section{Concepts}

\subsection{Survivorship Bias}

\begin{definition}[Survivorship Bias]
    The error of only including assets that \emph{survived} to the present date in
    a historical analysis.  Delisted stocks (bankruptcies, acquisitions,
    regulatory removals) are excluded, biasing aggregate returns upward.

    A backtest universe must be constructed \textbf{as-of} each historical date,
    including stocks that existed at that date even if they were later delisted.
\end{definition}

\subsection{Liquidity Filters}

Common filters for a ``tradeable'' universe:

\begin{itemize}
    \item \textbf{Minimum price}: Exclude stocks below \$1--\$5 (penny stock noise).
    \item \textbf{Minimum volume}: Exclude stocks with average daily volume below
          a threshold (e.g., 100{,}000 shares) --- you can't trade what you can't
          fill.
    \item \textbf{Market cap}: Exclude micro-caps if your strategy can't handle
          the illiquidity and volatility.
    \item \textbf{Listing exchange}: Exclude OTC/pink-sheet stocks for most
          strategies.
\end{itemize}

\begin{remark}
    Every filter is a tradeoff.  Excluding small-caps removes noise but also
    removes potential alpha.  The filter parameters should be \emph{explicit and
    documented}, not buried in a data loader.
\end{remark}

\subsection{Sector and Industry Mapping}

Cross-sectional analysis (Part~III) requires grouping stocks by sector.  Standard
classifications:

\begin{itemize}
    \item \textbf{GICS} (Global Industry Classification Standard): 11 sectors,
          24 industry groups --- the institutional standard.
    \item \textbf{SIC / NAICS}: Older classification systems, still used in SEC
          filings.
\end{itemize}

Sector membership is time-varying (companies change sectors after M\&A or
business pivots).  A proper implementation uses point-in-time sector assignments.

% ── Implementation ────────────────────────────────────────────────────────────
\section{Implementation Sketch}

A universe construction module should provide:

\begin{lstlisting}
def build_universe(
    as_of_date: date,
    min_price: float = 1.0,
    min_avg_volume: int = 100_000,
    exclude_otc: bool = True,
) -> list[str]:
    """
    Return a list of tradeable tickers as-of a given date.
    Point-in-time: includes stocks that were active on as_of_date,
    even if they were later delisted.
    """
\end{lstlisting}

The key constraint is \textbf{point-in-time correctness}: the function must not
use information from after \texttt{as\_of\_date}.

\begin{labmod}
    \textbf{Build it}: Implement a minimal universe filter.  Compare the number
    of tickers in the filtered universe vs.\ the full ticker list at several
    historical dates.  How much does the universe size change over time?
\end{labmod}

% ── Experiment ────────────────────────────────────────────────────────────────
\section{Experiment: Survivorship Bias Quantification}

\begin{experiment}
    \textbf{Goal}: Measure the impact of survivorship bias on a simple momentum
    strategy.

    \textbf{Method}:
    \begin{enumerate}
        \item Run a 12-month momentum strategy on today's S\&P~500 constituents
              (survivorship-biased universe).
        \item Run the same strategy on the point-in-time universe (historical
              constituents at each rebalance date).
        \item Compare CAGR and Sharpe ratio.
    \end{enumerate}

    \textbf{Expected}: The biased backtest will show higher returns (failed
    companies excluded).  The difference is the ``survivorship premium.''

    \textbf{Reflection space}: \emph{(fill in after running)}
\end{experiment}


% ══════════════════════════════════════════════════════════════════════════════
% PART II: STRATEGY & BACKTESTING
% ══════════════════════════════════════════════════════════════════════════════

\part{Strategy and Backtesting}
\label{part:backtest}

\begin{quote}
\itshape
A backtest is a simulation of what \emph{would have happened} if you had traded a
strategy in the past.  It is the primary tool for evaluating ideas --- and the
primary source of self-deception if done carelessly.  This part covers how to
design strategies, measure performance, and avoid the biases that make bad
strategies look good.
\end{quote}


% ──────────────────────────────────────────────────────────────────────────────
\chapter{Technical Indicators and Signal Design}
\label{ch:indicators}

% ── Motivation ────────────────────────────────────────────────────────────────
\section{Motivation}

An indicator transforms raw price/volume data into a derived quantity that is
(hopefully) easier to make trading decisions from.  Every rule-based strategy
starts with indicators; every factor-based strategy (Part~III) is built from
them.

The engineering question is not ``which indicator is best'' --- it is ``how do I
organize indicators so they are composable, testable, and reusable across
strategies?''

\textbf{Interview relevance}: You won't be asked to recite indicator formulas.
You \emph{will} be asked: ``How is your indicator pipeline structured?  Can you
add a new indicator without touching the strategy code?''

% ── Concepts ──────────────────────────────────────────────────────────────────
\section{Concepts}

\subsection{Indicator Taxonomy}

\begin{itemize}
    \item \textbf{Trend}: Moving averages (SMA, EMA, BBI), MACD, ADX ---
          capture the direction and strength of price movement.
    \item \textbf{Mean reversion}: Bollinger Bands, RSI, z-score of price
          relative to a moving average --- capture deviation from ``normal.''
    \item \textbf{Volume}: OBV, VWAP, Accumulation/Distribution --- capture the
          conviction behind price moves.
    \item \textbf{Volatility}: ATR, Bollinger bandwidth, realized volatility ---
          capture the magnitude of price fluctuations.
\end{itemize}

\subsection{The Registry Pattern}

A clean indicator system uses a \textbf{registry}: each indicator registers
itself via a decorator, and the strategy code requests indicators by name without
knowing their implementation details.

\begin{lstlisting}
@register_indicator("bbiboll")
def compute_bbiboll(df: pl.DataFrame) -> pl.DataFrame:
    """Compute BBI + Bollinger deviation percentile rank."""
    ...
\end{lstlisting}

Benefits:
\begin{itemize}
    \item Adding a new indicator requires zero changes to existing code.
    \item Indicators can declare whether they need per-ticker grouping.
    \item The registry enables bulk computation: ``apply all registered
          indicators to this DataFrame.''
\end{itemize}

% ── Implementation ────────────────────────────────────────────────────────────
\section{Implementation Sketch}

The indicator subsystem has three layers:

\begin{enumerate}
    \item \textbf{Registry} --- a global dictionary mapping names to
          functions, populated by decorator registration.
    \item \textbf{Base applicator} --- a helper that applies a registered
          indicator function over grouped data (group-by-ticker).
    \item \textbf{Indicator modules} --- one file per indicator (or family),
          each file auto-imported to trigger registration.
\end{enumerate}

\begin{labmod}
    \textbf{Build it}: Implement a new indicator (e.g., RSI or ATR) using the
    registry pattern.  Verify it works by applying it to a single ticker's OHLCV
    DataFrame, then to the full universe via the group-by applicator.  Confirm
    that no strategy code needed changes.
\end{labmod}


% ──────────────────────────────────────────────────────────────────────────────
\chapter{Rule-Based Strategies}
\label{ch:rule-strategy}

% ── Motivation ────────────────────────────────────────────────────────────────
\section{Motivation}

Before building factor models or ML pipelines, every quant should build at least
one rule-based strategy end to end: from signal definition through execution logic
to PnL calculation.  This exercises the full stack and reveals problems
(data bugs, look-ahead leaks, position tracking errors) that are invisible in
theory.

% ── Concepts ──────────────────────────────────────────────────────────────────
\section{Concepts}

\subsection{Strategy as a State Machine}

A rule-based strategy is a \textbf{state machine} with states
\{Flat, Long, Short\} and transitions triggered by indicator conditions:

\begin{itemize}
    \item \textbf{Entry signal}: Condition that triggers opening a position
          (e.g., BBI deviation crosses below the 10th percentile $\Rightarrow$
          buy).
    \item \textbf{Exit signal}: Condition that triggers closing
          (e.g., deviation returns above the 50th percentile, or stop-loss
          hit).
    \item \textbf{Position sizing}: How much capital to allocate per trade
          (equal weight, volatility-targeted, Kelly --- see Chapter~\ref{ch:portfolio}).
\end{itemize}

\subsection{The Strategy Base Class}

An abstract base class enforces a contract:

\begin{lstlisting}
class StrategyBase(ABC):
    @abstractmethod
    def generate_signals(self, data: pl.DataFrame) -> pl.DataFrame:
        """Return a DataFrame with a 'signal' column."""

    def run(self, data: pl.DataFrame) -> pl.DataFrame:
        """Apply signals, compute positions, return trades."""
\end{lstlisting}

Every concrete strategy inherits this interface.  The backtest engine only
depends on the base class --- it doesn't know whether the strategy uses BBIBOLL,
SMA crossover, or a neural network.

\subsection{Case Study: BBIBOLL}

The project's first complete strategy combines two ideas:

\begin{enumerate}
    \item \textbf{BBI (Bull and Bear Index)}: The average of four SMAs
          (3, 6, 12, 24 periods).  Smoother than any single MA.
    \item \textbf{Bollinger deviation}: The distance of price from BBI,
          expressed as a percentile rank within a rolling window.  Low
          percentile $=$ price is unusually far below the mean $=$ potential
          mean-reversion buy signal.
\end{enumerate}

Entry: deviation percentile $<$ threshold.  Exit: percentile recovers, or
stop-loss / take-profit triggers.

\begin{remark}
    BBIBOLL is a \textbf{time-series signal}: computed from each stock's own
    history.  In Part~III, we will reframe it as a \textbf{cross-sectional
    factor} by ranking the signal across all stocks at each date.  This
    reframing is the key conceptual leap from ``strategy'' to ``alpha.''
\end{remark}

\begin{labmod}
    \textbf{Build it}: Implement a minimal SMA-crossover strategy using
    the strategy base class interface.  Run it on one ticker for one year.
    Inspect the trades DataFrame: are entry/exit dates correct?  Does
    position tracking handle edge cases (e.g., signal on the last day)?
\end{labmod}


% ──────────────────────────────────────────────────────────────────────────────
\chapter{Backtest Engine and Performance Metrics}
\label{ch:engine}

% ── Motivation ────────────────────────────────────────────────────────────────
\section{Motivation}

The engine is the machinery that turns a strategy's signals into simulated
equity curves.  The performance metrics are the language you use to describe
whether the results are good, bad, or meaningless.

Both are mechanical --- the intellectual challenge is in the \emph{next} chapter
(bias and validation).  But you need the mechanics first.

% ── Concepts ──────────────────────────────────────────────────────────────────
\section{Concepts}

\subsection{The Backtest Loop}

At its core, a backtest iterates over time steps and applies:

\begin{enumerate}
    \item \textbf{Signal generation}: Run the strategy on data available up to
          time $t$.
    \item \textbf{Order creation}: Translate signals into hypothetical orders.
    \item \textbf{Execution simulation}: Fill orders at assumed prices
          (e.g., next-day open) with assumed costs.
    \item \textbf{Portfolio update}: Update positions, cash, and equity.
\end{enumerate}

\begin{pitfall}
    \textbf{Look-ahead in the loop.}  The most common backtest bug: using the
    close price at time $t$ to generate a signal and then assuming you can
    execute at that same close price.  In reality, you can only act on
    information \emph{after} observing it --- execution should happen at
    $t+1$ open at the earliest.
\end{pitfall}

\subsection{Return Metrics}

\begin{definition}[Sharpe Ratio]
    \[
        SR = \frac{\mathbb{E}[r_p - r_f]}{\sigma(r_p - r_f)}
    \]
    Annualized: $SR_{\text{annual}} = SR_{\text{daily}} \times \sqrt{252}$.

    The Sharpe ratio is the single most quoted performance metric.  It measures
    \textbf{risk-adjusted return}: how much return you earn per unit of
    volatility.  See Appendix~\ref{sec:math-stats} for the underlying
    statistics.
\end{definition}

\begin{definition}[Sortino Ratio]
    Like Sharpe, but only penalizes \textbf{downside} volatility:
    \[
        \text{Sortino} = \frac{\mathbb{E}[r_p - r_f]}{\sigma_{\text{downside}}}
    \]
    Preferred when returns are asymmetric (positive skew $\Rightarrow$ Sortino
    $>$ Sharpe).
\end{definition}

\begin{definition}[Maximum Drawdown]
    \[
        MDD = \max_{t} \left( \frac{\text{Peak}_t - \text{Trough}_t}{\text{Peak}_t} \right)
    \]
    The largest peak-to-trough decline.  For a live trader, MDD is the metric
    that determines whether you get fired or your fund gets redeemed.
\end{definition}

\begin{definition}[CAGR and Calmar Ratio]
    \[
        CAGR = \left( \frac{V_T}{V_0} \right)^{1/T} - 1
        \qquad\qquad
        \text{Calmar} = \frac{CAGR}{|MDD|}
    \]
\end{definition}

\subsection{Trade-Level Metrics}

\begin{definition}[Win Rate and Payoff Ratio]
    \begin{align}
        \text{Win Rate} &= \frac{\text{winning trades}}{\text{total trades}} \\[4pt]
        \text{Payoff Ratio} &= \frac{\text{average win}}{\text{average loss}}
    \end{align}
    A strategy can be profitable with win rate $< 50\%$ if the payoff ratio is
    high (trend following), or with low payoff ratio if win rate is high (mean
    reversion).
\end{definition}

% ── Implementation ────────────────────────────────────────────────────────────
\section{Implementation Sketch}

The engine consists of three decoupled components:

\begin{enumerate}
    \item \textbf{Runner}: Wires strategy $+$ data $+$ analyzer; produces the
          equity curve.
    \item \textbf{PerformanceAnalyzer}: Computes all metrics from the equity
          curve and trade log.
    \item \textbf{Visualizer}: Plots equity curves, drawdowns, monthly
          heatmaps, and candlestick charts with signal overlays.
\end{enumerate}

\begin{labmod}
    \textbf{Build it}: Run the BBIBOLL strategy on a rolling set of weekly
    start dates (60+ windows).  Collect the Sharpe ratio from each window.
    Plot the distribution of Sharpe ratios.  Is the strategy consistently
    positive, or does a few good windows dominate?
\end{labmod}


% ──────────────────────────────────────────────────────────────────────────────
\chapter{Backtest Bias and Validation}
\label{ch:bias}

% ── Motivation ────────────────────────────────────────────────────────────────
\section{Motivation}

\textbf{This is the most important chapter in Part~II.}  A backtest that looks
good but contains biases is worse than no backtest at all --- it gives false
confidence that costs real money.

Every quant interview will probe your understanding of backtest biases.  Saying
``my backtest Sharpe is 2.5'' without being able to explain \emph{why it might be
wrong} is a red flag.

% ── Concepts ──────────────────────────────────────────────────────────────────
\section{Concepts}

\subsection{The Three Cardinal Biases}

\begin{definition}[Look-Ahead Bias]
    Using information at time $t$ that was not actually available at time $t$.

    Examples:
    \begin{itemize}
        \item Using adjusted prices that incorporate \emph{future} splits.
        \item Using the full-sample mean to normalize signals.
        \item Generating a signal at close and executing at the same close.
    \end{itemize}
\end{definition}

\begin{definition}[Survivorship Bias]
    Only backtesting on assets that survived to the present day.  Delisted
    stocks (often the worst performers) are excluded, biasing returns upward.

    See Chapter~\ref{ch:universe} for how universe construction addresses this.
\end{definition}

\begin{definition}[Data Snooping / Overfitting]
    Testing many parameter combinations on the same dataset and selecting the
    best result.  The more combinations you test, the more likely one will
    look good \emph{by chance}.

    \textbf{Deflated Sharpe Ratio} (Bailey \& L\'{o}pez de Prado, 2014):
    adjusts the Sharpe ratio for the number of strategies tested.
\end{definition}

\subsection{Walk-Forward Validation}

\begin{definition}[Walk-Forward Optimization (WFO)]
    \begin{enumerate}
        \item Split data into rolling windows: $[\text{train} \mid \text{test}]$.
        \item Optimize parameters on the train window.
        \item Evaluate on the test window (\textbf{no re-optimization}).
        \item Slide the window forward and repeat.
        \item Concatenate all test-window results for the ``true'' out-of-sample
              performance.
    \end{enumerate}
\end{definition}

\begin{remark}
    The project's rolling weekly backtest (60+ date windows with fixed
    parameters) resembles walk-forward testing but is \emph{not} true WFO
    because parameters are not re-optimized inside each training window.
    Upgrading to true WFO is a planned enhancement.
\end{remark}

\subsection{Statistical Significance of Backtest Results}

\begin{definition}[$t$-Statistic of the Sharpe Ratio]
    Is the backtest Sharpe distinguishable from zero?
    \[
        t = SR \times \sqrt{n}
    \]
    where $n$ is the number of return periods.  A rule of thumb: $|t| > 2$
    (equivalently, $p < 0.05$) suggests the result is unlikely due to chance
    alone.

    See Appendix~\ref{sec:math-stats} for the underlying hypothesis testing
    framework.
\end{definition}

\begin{pitfall}
    A statistically significant Sharpe does not mean an \textbf{economically}
    significant strategy.  A Sharpe of 0.3 with $|t| = 2.5$ (from 20 years of
    data) is ``real'' but may not cover transaction costs.
\end{pitfall}

% ── Experiment ────────────────────────────────────────────────────────────────
\section{Experiment: How Fragile Is the BBIBOLL Sharpe?}

\begin{experiment}
    \textbf{Goal}: Stress-test the BBIBOLL strategy's reported Sharpe ratio.

    \textbf{Method}:
    \begin{enumerate}
        \item Compute the Sharpe from the full sample.
        \item Compute the $t$-statistic.  Is $|t| > 2$?
        \item Remove the best 5 trading days.  How much does the Sharpe drop?
        \item Run on a different time period (out-of-sample).  Does the Sharpe
              survive?
        \item Shuffle trade entry dates randomly (permutation test).  What
              Sharpe do you get from pure randomness?
    \end{enumerate}

    \textbf{Reflection space}: \emph{(fill in after running)}
\end{experiment}


% ══════════════════════════════════════════════════════════════════════════════
% PART III: ALPHA RESEARCH
% ══════════════════════════════════════════════════════════════════════════════

\part{Alpha Research}
\label{part:alpha}

\begin{quote}
\itshape
This is the inflection point.  Parts~I and~II make you a ``trader who codes.''
Part~III makes you a ``quant researcher.''  The core shift: from asking ``does
this strategy make money?'' to asking ``does this signal predict returns, and
how do I measure that rigorously?''
\end{quote}


% ──────────────────────────────────────────────────────────────────────────────
\chapter{From Rules to Factors}
\label{ch:rules-to-factors}

% ── Motivation ────────────────────────────────────────────────────────────────
\section{Motivation}

In Part~II you built a rule-based strategy: ``if deviation percentile $< 10$,
buy.''  This works, but it has a problem: the threshold is arbitrary, the
evaluation is binary (profitable or not), and there is no principled way to
combine it with other ideas.

The factor framework solves this.  Instead of asking ``should I buy?'', you ask:
``across all 5{,}000 stocks today, does a higher deviation percentile predict
higher future returns?''

This reframing unlocks:
\begin{itemize}
    \item \textbf{Continuous evaluation}: IC measures the strength of the signal,
          not just ``profitable or not.''
    \item \textbf{Composability}: Multiple factors can be combined into a single
          portfolio signal.
    \item \textbf{Decomposition}: You can attribute portfolio returns to
          individual factors.
\end{itemize}

\textbf{Interview relevance}: This chapter covers the exact vocabulary of quant
researcher interviews.  ``What is the IC of your factor?'' is the first question.

% ── Concepts ──────────────────────────────────────────────────────────────────
\section{Concepts}

\subsection{Time-Series vs.\ Cross-Sectional Thinking}

\begin{definition}[Time-Series Signal]
    A signal computed from a single asset's own history.  Example: 20-day
    momentum return for AAPL, BBIBOLL deviation for NVDA\@.  The signal value
    for AAPL does not depend on any other stock.
\end{definition}

\begin{definition}[Cross-Sectional Signal (Factor)]
    A signal that \textbf{ranks or normalizes} values across the entire stock
    universe at each point in time.  Example: the z-score of PE ratio across
    all stocks on a given date.

    The key property: the factor value for stock $i$ on date $t$ depends on
    \emph{all other stocks} on that date.
\end{definition}

The BBIBOLL deviation percentile is a time-series signal.  To use it
cross-sectionally:
\[
    \text{factor}_{i,t} = \frac{\text{deviation}_{i,t} - \overline{\text{deviation}}_t}
                               {\sigma(\text{deviation}_t)}
\]
where the mean and standard deviation are computed across all stocks $i$ at
date $t$.  This produces a z-scored cross-sectional factor.

\subsection{The Alpha Pipeline}

The standard factor research workflow:

\begin{enumerate}
    \item \textbf{Signal construction}: Compute a raw signal for each stock at
          each date.
    \item \textbf{Normalization}: Z-score or rank-transform cross-sectionally.
    \item \textbf{Evaluation}: Compute IC, IR, decay, turnover.
    \item \textbf{Refinement}: Neutralize for sector, size, or other known
          factors.
    \item \textbf{Combination}: Blend with other factors into a composite
          signal.
    \item \textbf{Portfolio construction}: Map the composite signal to
          positions (Part~IV).
\end{enumerate}

Each step is a separate module.  This chapter covers steps 1--2; the next three
chapters cover steps 3--5.

\begin{labmod}
    \textbf{Build it}: Take the BBIBOLL deviation signal (a time-series signal)
    and convert it to a cross-sectional factor.  At each date, z-score the
    deviation across all stocks in the liquid universe (Chapter~\ref{ch:universe}).
    Plot the distribution of factor values on a single date --- is it
    approximately normal?
\end{labmod}


% ──────────────────────────────────────────────────────────────────────────────
\chapter{Factor Evaluation}
\label{ch:factor-eval}

% ── Motivation ────────────────────────────────────────────────────────────────
\section{Motivation}

You've built a factor.  Is it any good?  ``Running a backtest'' is not the
answer --- backtests conflate signal quality with portfolio construction,
execution, and luck.  Factor evaluation isolates the question: \textbf{does
this signal predict future returns?}

This chapter covers the four metrics that every quant uses to answer that
question.

% ── Concepts ──────────────────────────────────────────────────────────────────
\section{Concepts}

\subsection{Information Coefficient (IC)}

\begin{definition}[IC]
    The cross-sectional rank correlation (Spearman's $\rho$) between a factor's
    signal values and subsequent forward returns:
    \[
        IC_t = \rho_{\text{rank}}\!\left(\text{signal}_{t},\; r_{t \to t+h}\right)
    \]
    where the correlation is computed across all stocks in the universe at
    date~$t$, and $h$ is the forecast horizon (e.g., 5 trading days).
\end{definition}

\begin{remark}
    Why Spearman (rank) and not Pearson?
    \begin{itemize}
        \item Robust to outliers (extreme returns, extreme signal values).
        \item Does not assume a linear relationship.
        \item Invariant to monotonic transformations of the signal.
    \end{itemize}
    In practice, you should compute \emph{both} and compare.  A large gap
    between Pearson and Spearman IC suggests outlier sensitivity.
\end{remark}

\subsection{Information Ratio (IR)}

\begin{definition}[IR]
    \[
        IR = \frac{\overline{IC}}{\sigma(IC)}
    \]
    where $\overline{IC}$ is the time-series mean of $IC_t$ and $\sigma(IC)$ is
    its standard deviation.

    IR measures the \textbf{consistency} of predictive power.  A factor with
    $|\overline{IC}| = 0.03$ and $IR = 0.5$ is far more valuable than one with
    $|\overline{IC}| = 0.05$ and $IR = 0.2$.  Consistency beats magnitude.
\end{definition}

\begin{remark}
    Rule of thumb (Qian, Hua \& Sorensen):
    \begin{center}
    \begin{tabular}{@{}ll@{}}
    \toprule
    $|IR|$ & Interpretation \\
    \midrule
    $> 0.5$       & Strong, tradeable factor \\
    $0.3$--$0.5$  & Decent, worth including in a composite \\
    $< 0.3$       & Weak or unreliable \\
    \bottomrule
    \end{tabular}
    \end{center}
\end{remark}

\subsection{IC Decay}

\begin{definition}[IC Decay Curve]
    Plot $\overline{IC}(h)$ for increasing forecast horizons
    $h = 1, 2, 5, 10, 20$ days.  The shape reveals the signal's natural
    trading frequency:
    \begin{itemize}
        \item \textbf{Fast decay} (IC drops sharply after 1--2 days):
              short-term signal, requires frequent trading.
        \item \textbf{Slow decay} (IC persists for weeks):
              medium-frequency signal, lower turnover.
        \item \textbf{Non-monotonic decay}: possible look-ahead bias or
              seasonal effects --- investigate.
    \end{itemize}
\end{definition}

\begin{definition}[Half-Life]
    The number of periods until the IC drops to half its peak value.
    Approximate from the AR(1) coefficient $\phi$ of the IC time series:
    \[
        t_{1/2} = -\frac{\ln 2}{\ln |\phi|}
    \]
\end{definition}

\subsection{Turnover}

\begin{definition}[Factor Turnover]
    How much the factor's ranking changes from one period to the next:
    \[
        \text{Turnover}_t = \frac{1}{N} \sum_{i=1}^{N}
            |w_{i,t} - w_{i,t-1}|
    \]
    where $w_{i,t}$ is the normalized factor weight for stock $i$ at time $t$.

    High turnover $\Rightarrow$ high transaction costs $\Rightarrow$ the factor
    needs a higher gross IC to be profitable after costs.
\end{definition}

\begin{definition}[Capacity]
    The maximum AUM a factor can support before market impact erodes its alpha.
    Small-cap momentum: high IC, low capacity.  Large-cap value: lower IC,
    high capacity.  Capacity is a function of the liquidity of the stocks the
    factor trades.
\end{definition}

% ── Implementation ────────────────────────────────────────────────────────────
\section{Implementation Sketch}

A factor evaluation module should provide:

\begin{lstlisting}
class FactorAnalyzer:
    def __init__(self, signal: pl.DataFrame, returns: pl.DataFrame):
        """
        signal: (date, ticker, signal_value)
        returns: (date, ticker, forward_return_1d, ..., forward_return_20d)
        """

    def ic_series(self, horizon: int = 5) -> pl.DataFrame:
        """Cross-sectional Spearman IC at each date."""

    def ir(self, horizon: int = 5) -> float:
        """Mean IC / Std IC."""

    def ic_decay(self, horizons: list[int]) -> pl.DataFrame:
        """Mean IC at each horizon."""

    def turnover(self) -> pl.DataFrame:
        """Daily factor turnover."""
\end{lstlisting}

\begin{labmod}
    \textbf{Build it}: Implement \texttt{FactorAnalyzer} with at least
    \texttt{ic\_series()} and \texttt{ir()}.  Run it on the cross-sectional
    BBIBOLL factor from Chapter~\ref{ch:rules-to-factors}.  Report the mean IC,
    IR, and $t$-statistic (is $|t| > 2$?).
\end{labmod}

% ── Experiment ────────────────────────────────────────────────────────────────
\section{Experiment: BBIBOLL Factor Evaluation}

\begin{experiment}
    \textbf{Goal}: Evaluate the BBIBOLL deviation signal as a cross-sectional
    factor.

    \textbf{Method}:
    \begin{enumerate}
        \item Compute the cross-sectional z-scored BBIBOLL deviation factor for
              the liquid universe.
        \item Compute the 5-day forward return for each stock at each date.
        \item Run \texttt{FactorAnalyzer}: compute IC time series, mean IC, IR,
              $t$-stat.
        \item Plot the IC decay curve for horizons 1, 2, 5, 10, 20 days.
        \item Compute daily turnover.
    \end{enumerate}

    \textbf{Questions to answer}:
    \begin{itemize}
        \item Is the mean IC positive or negative?  (Directional prediction.)
        \item Is $|t| > 2$?  (Statistically significant?)
        \item What is the natural horizon?  (Where IC peaks before decaying.)
        \item Is turnover manageable?  (Can you trade this weekly?)
    \end{itemize}

    \textbf{Reflection space}: \emph{(fill in after running)}
\end{experiment}


% ──────────────────────────────────────────────────────────────────────────────
\chapter{Factor Construction}
\label{ch:factor-construction}

% ── Motivation ────────────────────────────────────────────────────────────────
\section{Motivation}

A raw signal is rarely optimal.  Cross-sectional normalization, sector
neutralization, and outlier handling can substantially improve IC and IR ---
or reveal that the signal was just a proxy for a well-known risk factor.

This chapter covers the standard transformations every factor goes through
before it is ready for combination and portfolio construction.

% ── Concepts ──────────────────────────────────────────────────────────────────
\section{Concepts}

\subsection{Normalization}

\begin{definition}[Z-Score Normalization]
    \[
        z_{i,t} = \frac{x_{i,t} - \bar{x}_t}{\sigma_{x,t}}
    \]
    Applied cross-sectionally: ensures zero mean, unit variance across the
    stock universe at each date.
\end{definition}

\begin{definition}[Rank Normalization]
    Replace raw values with their cross-sectional rank, then map to $[-1, 1]$
    or to a standard normal via the inverse CDF\@.  More robust to outliers
    than z-scoring.
\end{definition}

\begin{definition}[Winsorization]
    Cap extreme values at a percentile threshold (e.g., 1st and 99th
    percentiles).  Applied \emph{before} z-scoring to prevent outliers from
    dominating the mean and standard deviation.
\end{definition}

\subsection{Neutralization}

\begin{definition}[Sector Neutralization]
    Demean the factor within each sector:
    \[
        z_{i,t}^{\text{neutral}} = z_{i,t} - \bar{z}_{\text{sector}(i),t}
    \]
    Purpose: remove sector-level effects so the factor captures stock-specific
    alpha, not sector rotation.
\end{definition}

\begin{definition}[Size and Beta Neutralization]
    Regress the factor on log-market-cap and/or market beta; use the residual:
    \[
        z^{\perp} = z - \hat{z}
    \]
    where $\hat{z}$ is the fitted value from the regression.  This isolates
    alpha from well-known risk premiums.
\end{definition}

\subsection{Orthogonalization}

\begin{definition}[Factor Orthogonalization]
    Given existing factors $F_1, \ldots, F_k$, orthogonalize a new factor
    $F_{k+1}$ by regressing it on the existing set and taking the residual:
    \[
        F_{k+1}^{\perp} = F_{k+1} - X(X^TX)^{-1}X^T F_{k+1}
    \]
    where $X = [F_1 \mid \cdots \mid F_k]$.

    The residual $F_{k+1}^{\perp}$ contains only the information in $F_{k+1}$
    that is \textbf{not} already captured by the existing factors.
\end{definition}

\begin{remark}
    Orthogonalization order matters.  $F_2$ orthogonalized against $F_1$ is
    different from $F_1$ orthogonalized against $F_2$.  The first factor in
    the sequence ``wins'' more of the shared variance.
\end{remark}

% ── Implementation ────────────────────────────────────────────────────────────
\section{Implementation Sketch}

A factor preprocessing pipeline:

\begin{lstlisting}
def preprocess_factor(
    raw_signal: pl.DataFrame,       # (date, ticker, value)
    sectors: pl.DataFrame,          # (ticker, sector)
    market_cap: pl.DataFrame,       # (date, ticker, log_mcap)
    winsorize_pct: float = 0.01,
    neutralize: list[str] = ["sector"],
) -> pl.DataFrame:
    """
    1. Winsorize at (pct, 1-pct) cross-sectionally.
    2. Z-score cross-sectionally.
    3. Neutralize for specified factors (sector, size).
    4. Return the preprocessed factor DataFrame.
    """
\end{lstlisting}

\begin{labmod}
    \textbf{Build it}: Implement \texttt{preprocess\_factor()} with at least
    z-scoring and sector neutralization.  Compare the IC of the raw BBIBOLL
    factor vs.\ the sector-neutralized version.  Does neutralization improve
    IR?
\end{labmod}

% ── Experiment ────────────────────────────────────────────────────────────────
\section{Experiment: Neutralization Impact}

\begin{experiment}
    \textbf{Goal}: Measure how sector neutralization affects the BBIBOLL
    factor.

    \textbf{Method}:
    \begin{enumerate}
        \item Compute the raw z-scored BBIBOLL factor.
        \item Apply sector neutralization using GICS sector assignments.
        \item Compute IC, IR, and turnover for both versions.
        \item Compare.
    \end{enumerate}

    \textbf{Hypothesis}: If BBIBOLL is mostly capturing sector rotation
    (e.g., ``tech is cheap relative to its own history''), neutralization will
    \emph{destroy} the IC\@.  If it captures stock-specific mean reversion,
    neutralization will \emph{preserve or improve} it.

    \textbf{Reflection space}: \emph{(fill in after running)}
\end{experiment}


% ──────────────────────────────────────────────────────────────────────────────
\chapter{Factor Combination}
\label{ch:factor-combo}

% ── Motivation ────────────────────────────────────────────────────────────────
\section{Motivation}

No single factor is strong enough to build a portfolio on.  The power of
systematic investing comes from \textbf{combining} many weak signals into a
composite that is more stable than any individual factor.

This chapter covers the main combination methods, from the simplest
(equal-weight) to the most sophisticated (mean-variance on IC covariance).

% ── Concepts ──────────────────────────────────────────────────────────────────
\section{Concepts}

\subsection{Equal-Weight and IC-Weight}

The simplest composite: $\text{composite}_{i,t} = \sum_k w_k \cdot f_{k,i,t}$
where $f_k$ are z-scored factors.

\begin{itemize}
    \item \textbf{Equal-weight}: $w_k = 1/K$ for all factors.  No estimation
          risk, but ignores differences in signal quality.
    \item \textbf{IC-weight}: $w_k \propto \overline{IC}_k$.  Allocates more
          weight to stronger factors.  Requires reliable IC estimates.
\end{itemize}

\subsection{Mean-Variance on Factor ICs}

\begin{definition}[Factor Mean-Variance Combination]
    Treat each factor's IC time series as a ``return stream.''  Apply
    mean-variance optimization:
    \[
        w^* = \arg\max_w \left( w^T \mu_{IC} - \frac{\lambda}{2}
              w^T \Sigma_{IC}\, w \right)
    \]
    where $\mu_{IC}$ is the vector of mean ICs and $\Sigma_{IC}$ is the
    covariance matrix of IC time series.

    This accounts for both factor quality (mean IC) and factor correlation
    (diversification benefit from uncorrelated factors).
\end{definition}

\subsection{Risk Parity}

\begin{definition}[Risk Parity Combination]
    Weight factors so that each contributes equally to the composite signal's
    variance:
    \[
        w_k \propto \frac{1}{\sigma_{IC,k}}
    \]
    Advantage: does not require estimating expected IC (only volatility), so it
    is more robust to estimation error.
\end{definition}

\begin{remark}
    In practice, start with equal-weight.  Only graduate to IC-weight or
    mean-variance when you have enough factors ($\geq 5$) and enough history
    ($\geq 2$ years of daily IC) to estimate the inputs reliably.
\end{remark}

\begin{labmod}
    \textbf{Build it}: Once you have $\geq 2$ factors (e.g., BBIBOLL deviation
    + 20-day momentum), implement equal-weight and IC-weight combination.
    Compare the composite's IC and IR against each individual factor.  Does
    combination improve IR?
\end{labmod}

% ── Experiment ────────────────────────────────────────────────────────────────
\section{Experiment: Two-Factor Composite}

\begin{experiment}
    \textbf{Goal}: Build and evaluate a two-factor composite signal.

    \textbf{Method}:
    \begin{enumerate}
        \item Factor A: Cross-sectional z-scored BBIBOLL deviation.
        \item Factor B: 20-day momentum (cross-sectional z-scored).
        \item Compute the Spearman correlation between A and B (across stocks,
              averaged over time).  Are they diversifying?
        \item Form composites: equal-weight and IC-weight.
        \item Evaluate each with \texttt{FactorAnalyzer}: IC, IR, decay.
    \end{enumerate}

    \textbf{Expected}: If A and B are negatively correlated (mean-reversion vs.\
    momentum), the composite's IR should be higher than either alone.

    \textbf{Reflection space}: \emph{(fill in after running)}
\end{experiment}


% ══════════════════════════════════════════════════════════════════════════════
% PART IV: PORTFOLIO, RISK & EXECUTION
% ══════════════════════════════════════════════════════════════════════════════

\part{Portfolio, Risk, and Execution}
\label{part:portfolio}

\begin{quote}
\itshape
A signal is not a portfolio.  Turning alpha signals into actual positions
requires risk management, position sizing, and realistic execution cost
modeling.  This part covers the bridge from ``I have a good factor'' to
``I have a tradeable strategy.''
\end{quote}


% ──────────────────────────────────────────────────────────────────────────────
\chapter{Risk Measures}
\label{ch:risk}

% ── Motivation ────────────────────────────────────────────────────────────────
\section{Motivation}

Risk management answers the question: ``how bad can it get?''  Without it, a
strategy with positive expected return can still bankrupt you via a single
tail event.

\textbf{Interview relevance}: ``What risk measures do you use?'' and
``Explain the difference between VaR and CVaR'' are standard questions.

% ── Concepts ──────────────────────────────────────────────────────────────────
\section{Concepts}

\subsection{Value at Risk}

\begin{definition}[VaR]
    The maximum loss over a given time horizon at a given confidence level:
    \[
        P(L > \text{VaR}_\alpha) = 1 - \alpha
    \]
    Example: ``1-day 95\% VaR $= \$10{,}000$'' means there is a 5\% chance
    of losing more than \$10{,}000 in one day.

    Three estimation methods:
    \begin{itemize}
        \item \textbf{Historical}: Use the empirical return distribution.
        \item \textbf{Parametric}: Assume a distribution (e.g., Gaussian)
              and compute the quantile analytically.
        \item \textbf{Monte Carlo}: Simulate many return paths and compute
              the empirical quantile.
    \end{itemize}
\end{definition}

\begin{definition}[CVaR (Expected Shortfall)]
    \[
        \text{CVaR}_\alpha = \mathbb{E}[L \mid L > \text{VaR}_\alpha]
    \]
    The expected loss \emph{given that} the loss exceeds VaR\@.  CVaR captures
    tail risk better than VaR because it asks ``when things go bad, \emph{how}
    bad?''
\end{definition}

\subsection{Distribution Properties of Returns}

\begin{definition}[Skewness and Kurtosis]
    \begin{itemize}
        \item \textbf{Skewness}: $\gamma_1 = \mathbb{E}\!\left[\left(\frac{X-\mu}{\sigma}\right)^{\!3}\right]$.
              Negative skew $=$ more extreme losses than gains.
        \item \textbf{Excess Kurtosis}: $\kappa = \mathbb{E}\!\left[\left(\frac{X-\mu}{\sigma}\right)^{\!4}\right] - 3$.
              Positive $\kappa$ $=$ fat tails.
    \end{itemize}
    Equity returns typically have \textbf{negative skewness} and \textbf{positive
    excess kurtosis}.  Gaussian VaR underestimates tail risk.
\end{definition}

\begin{labmod}
    \textbf{Build it}: Compute the daily return distribution of the BBIBOLL
    strategy.  Plot the histogram against a fitted Gaussian.  Compute skewness,
    excess kurtosis, 95\% VaR (historical and parametric), and CVaR\@.  How much
    does the parametric VaR underestimate tail risk?
\end{labmod}


% ──────────────────────────────────────────────────────────────────────────────
\chapter{Portfolio Construction}
\label{ch:portfolio}

% ── Motivation ────────────────────────────────────────────────────────────────
\section{Motivation}

A factor tells you \emph{which} stocks to favor.  Portfolio construction tells
you \emph{how much} of each stock to hold.  The mapping from signal to position
is where risk management and alpha research meet.

% ── Concepts ──────────────────────────────────────────────────────────────────
\section{Concepts}

\subsection{Factor Exposure and CAPM}

\begin{definition}[Alpha and Beta]
    From the CAPM regression: $r_i - r_f = \alpha_i + \beta_i (r_m - r_f) + \epsilon_i$
    \begin{itemize}
        \item $\beta_i$: sensitivity to the market.  $\beta = 1$ means the stock
              moves with the market.
        \item $\alpha_i$: excess return not explained by market movement.
              \textbf{This is what we seek.}
    \end{itemize}
\end{definition}

\begin{definition}[Market-Neutral Portfolio]
    A portfolio with net $\beta \approx 0$: long stocks with positive alpha,
    short stocks with negative alpha, weighted so that the net market exposure
    cancels.  Returns are (approximately) independent of market direction.
\end{definition}

\subsection{Position Sizing}

\begin{definition}[Equal-Weight]
    Allocate $1/N$ of capital to each position.  Simple, no estimation risk,
    but ignores differences in signal strength and risk.
\end{definition}

\begin{definition}[Volatility Targeting]
    Scale position sizes inversely with each stock's recent volatility:
    \[
        w_i \propto \frac{\sigma_{\text{target}}}{\hat{\sigma}_i}
    \]
    This ensures each position contributes roughly equal risk to the portfolio.
\end{definition}

\begin{definition}[Kelly Criterion]
    For a portfolio with expected excess return vector $\mu$ and covariance
    matrix $\Sigma$:
    \[
        f^* = \Sigma^{-1} \mu
    \]
    Full Kelly maximizes long-run growth but has extreme drawdowns.
    In practice, \textbf{half-Kelly} ($f^*/2$) or even quarter-Kelly is used.
\end{definition}

\begin{labmod}
    \textbf{Build it}: Implement volatility-targeted position sizing.  Compare
    equal-weight vs.\ vol-targeted performance for the BBIBOLL strategy.
    Measure: Sharpe, MDD, Calmar.
\end{labmod}


% ──────────────────────────────────────────────────────────────────────────────
\chapter{Execution and Transaction Costs}
\label{ch:execution}

% ── Motivation ────────────────────────────────────────────────────────────────
\section{Motivation}

A backtest without transaction costs is a fantasy.  The gap between ``paper
alpha'' and ``live alpha'' is almost entirely explained by execution costs:
commissions, bid-ask spread, slippage, and market impact.

This chapter covers the mechanics of execution and how to model costs in a
backtest.

\textbf{Interview relevance}: ``How do you model transaction costs in your
backtest?'' is a question that separates candidates who have thought about
real trading from those who haven't.

% ── Concepts ──────────────────────────────────────────────────────────────────
\section{Concepts}

\subsection{Cost Components}

\begin{definition}[Bid-Ask Spread]
    $\text{Spread} = P_{\text{ask}} - P_{\text{bid}}$.  Crossing the spread to
    execute a market order costs approximately half the spread per side.
\end{definition}

\begin{definition}[Slippage]
    The difference between the expected execution price and the actual fill.
    Caused by: order book depth, timing delays, stale quotes.
\end{definition}

\begin{definition}[Market Impact]
    The price movement caused by your own order:
    \begin{itemize}
        \item \textbf{Temporary impact}: Price displacement during execution;
              reverts afterward.
        \item \textbf{Permanent impact}: Information leakage that shifts the
              fair price.
    \end{itemize}
    Square-root model (empirical):
    $\text{Impact} \approx \sigma \cdot \sqrt{V_{\text{order}} / V_{\text{daily}}}$
\end{definition}

\subsection{Execution Algorithms}

\begin{definition}[VWAP and TWAP]
    Algorithms that slice large orders to reduce impact:
    \begin{itemize}
        \item \textbf{VWAP}: Execute proportional to historical intraday volume
              profile.
        \item \textbf{TWAP}: Execute equally over fixed time intervals.
    \end{itemize}
\end{definition}

\subsection{Cost Modeling in Backtests}

Three levels of sophistication:

\begin{enumerate}
    \item \textbf{Fixed cost}: Deduct a fixed amount per trade (e.g., \$0.005/share).
          Simple but unrealistic for large or illiquid orders.
    \item \textbf{Spread model}: Deduct half the typical bid-ask spread per side.
          Requires spread data (or estimates from Amihud illiquidity).
    \item \textbf{Impact model}: Use the square-root model above.  Requires
          volume data.  Most realistic, but adds complexity.
\end{enumerate}

\begin{pitfall}
    Ignoring costs is the most common source of ``paper alpha'' that disappears
    in live trading.  Even a simple fixed-cost model is infinitely better than
    zero-cost.
\end{pitfall}

\begin{labmod}
    \textbf{Build it}: Add a transaction cost model to the backtest engine.
    Start with level~1 (fixed cost), then upgrade to level~2 (spread model).
    Re-run the BBIBOLL backtest with costs.  How much does the Sharpe drop?
    At what cost level does the strategy become unprofitable?
\end{labmod}

% ── Experiment ────────────────────────────────────────────────────────────────
\section{Experiment: Cost Sensitivity Analysis}

\begin{experiment}
    \textbf{Goal}: Determine how sensitive the BBIBOLL strategy's performance
    is to transaction costs.

    \textbf{Method}:
    \begin{enumerate}
        \item Run the backtest with costs at 0, 1, 5, 10, 20 bps per side.
        \item Plot Sharpe ratio as a function of cost.
        \item Find the breakeven cost (where Sharpe $= 0$).
    \end{enumerate}

    \textbf{Interpretation}: If the breakeven cost is $< 5$ bps, the strategy
    is probably not tradeable in practice.  If $> 20$ bps, it is robust.

    \textbf{Reflection space}: \emph{(fill in after running)}
\end{experiment}


% ══════════════════════════════════════════════════════════════════════════════
% PART V: MACHINE LEARNING
% ══════════════════════════════════════════════════════════════════════════════

\part{Machine Learning in Quantitative Finance}
\label{part:ml}

\begin{quote}
\itshape
Machine learning in quant finance is not about model complexity.  It is about
\textbf{feature engineering}, \textbf{proper validation}, and \textbf{respecting
the time-series structure} of financial data.  If you get these three things
right, a simple model will beat a complex one.
\end{quote}


% ──────────────────────────────────────────────────────────────────────────────
\chapter{Feature Engineering for Finance}
\label{ch:features}

% ── Motivation ────────────────────────────────────────────────────────────────
\section{Motivation}

The gap between ``ML engineer'' and ``quant ML researcher'' is not modeling ---
it is features.  A LightGBM model with 50 carefully engineered financial
features will outperform a deep neural network fed raw OHLCV\@.

% ── Concepts ──────────────────────────────────────────────────────────────────
\section{Concepts}

\subsection{Feature Categories}

\begin{definition}[Financial Feature Taxonomy]
    \begin{itemize}
        \item \textbf{Momentum}: Returns over multiple horizons
              (1d, 5d, 20d, 60d, 252d).
        \item \textbf{Mean reversion}: Distance from moving averages
              (e.g., the BBIBOLL deviation signal).
        \item \textbf{Volume}: Relative volume, OBV,
              volume-price divergence, Amihud illiquidity.
        \item \textbf{Volatility}: Realized vol (from intraday data), GARCH
              residuals, vol-of-vol.
        \item \textbf{Fundamental}: PE, PB, earnings yield, dividend yield,
              revenue growth (requires fundamental data).
        \item \textbf{Microstructure}: Bid-ask spread statistics, trade
              imbalance, Kyle's lambda (requires tick data).
    \end{itemize}
\end{definition}

\subsection{Feature Processing}

\begin{itemize}
    \item \textbf{Cross-sectional normalization}: Z-score or rank-transform
          features across stocks at each date (same as factor preprocessing
          in Chapter~\ref{ch:factor-construction}).
    \item \textbf{Lag features}: Use the signal value from $t-1$, $t-2$, etc.\
          as additional inputs.  Captures mean-reversion or momentum in the
          signal itself.
    \item \textbf{Interaction features}: Ratios or products of base features
          (e.g., momentum $\times$ volume = ``confirmed momentum'').
\end{itemize}

\begin{pitfall}
    \textbf{Leakage through normalization.}  If you z-score using the full
    sample mean/std (including future dates), you introduce look-ahead bias.
    Always normalize using only data available up to date $t$.
\end{pitfall}

\begin{labmod}
    \textbf{Build it}: Construct a feature matrix with $\geq 10$ features
    (momentum at 3 horizons, BBIBOLL deviation, OBV z-score, realized vol,
    etc.) for the liquid universe.  Inspect the cross-sectional distribution
    of each feature --- any extreme outliers?  Apply winsorization and z-scoring.
\end{labmod}


% ──────────────────────────────────────────────────────────────────────────────
\chapter{Time-Series Validation}
\label{ch:ts-validation}

% ── Motivation ────────────────────────────────────────────────────────────────
\section{Motivation}

Standard $k$-fold cross-validation is \textbf{invalid} for time series because
it uses future data to predict the past.  This is the single most common ML
mistake in quant finance --- and the easiest way to produce a backtest that
looks incredible on paper and fails immediately in live trading.

% ── Concepts ──────────────────────────────────────────────────────────────────
\section{Concepts}

\subsection{Validation Schemes}

\begin{definition}[TimeSeriesSplit]
    \begin{itemize}
        \item \textbf{Expanding window}: Train on $[0, t]$, test on
              $[t+1, t+h]$.  Increase $t$ each iteration.
        \item \textbf{Rolling window}: Train on $[t-w, t]$, test on
              $[t+1, t+h]$.  Fixed training window width $w$.
        \item \textbf{Purged / Embargoed}: Remove $g$ days of data around
              the train/test boundary to prevent label leakage from
              overlapping forward-return windows.
    \end{itemize}
\end{definition}

\begin{definition}[Walk-Forward Training]
    The ML equivalent of walk-forward backtesting (Chapter~\ref{ch:bias}):
    retrain the model periodically (e.g., monthly) on the most recent $N$
    months, predict the next month, concatenate all out-of-sample predictions.
\end{definition}

\begin{pitfall}
    \textbf{Hyperparameter tuning on the test set.}  If you tune
    \texttt{max\_depth} or \texttt{learning\_rate} by looking at test-set
    performance, you are overfitting to the test set.  Use a nested validation
    scheme: train / validation / test.
\end{pitfall}

\begin{labmod}
    \textbf{Build it}: Implement a \texttt{PurgedWalkForwardCV} class that
    yields (train, test) index pairs with a configurable embargo gap.  Verify
    that no test date appears in any training fold.
\end{labmod}


% ──────────────────────────────────────────────────────────────────────────────
\chapter{Tree Models and Ensembles}
\label{ch:trees}

% ── Motivation ────────────────────────────────────────────────────────────────
\section{Motivation}

In quant ML, \textbf{gradient-boosted trees} (LightGBM, XGBoost) are the
workhorse.  They handle tabular data naturally, capture nonlinear interactions,
and provide built-in feature importance.  A well-tuned LightGBM with proper
walk-forward validation is a strong baseline that is hard to beat with more
complex models.

% ── Concepts ──────────────────────────────────────────────────────────────────
\section{Concepts}

\subsection{Why Trees for Finance?}

\begin{itemize}
    \item \textbf{Tabular data}: Financial features are structured columns,
          not images or text.  Trees excel on tabular data.
    \item \textbf{Nonlinear interactions}: The relationship between momentum
          and future return may depend on volatility.  Trees capture this
          automatically via splits.
    \item \textbf{Feature importance}: Built-in gain/split importance helps
          with feature selection and factor discovery.
    \item \textbf{Speed}: LightGBM trains on $10^6$ rows $\times$ 50 features
          in seconds.
\end{itemize}

\subsection{Ensemble Methods}

\begin{definition}[Bagging vs.\ Boosting]
    \begin{itemize}
        \item \textbf{Bagging} (Random Forest): Train many trees on bootstrap
              samples; average predictions.  Reduces \emph{variance}.
        \item \textbf{Boosting} (LightGBM, XGBoost): Train trees sequentially,
              each correcting the previous one's errors.  Reduces \emph{bias}.
    \end{itemize}
    In practice, boosting dominates for tabular prediction tasks.
\end{definition}

\subsection{Hyperparameter Discipline}

Key hyperparameters for LightGBM in a quant context:

\begin{itemize}
    \item \texttt{max\_depth}: Limit to 4--6 to prevent overfitting to noise.
    \item \texttt{num\_leaves}: Control tree complexity (related to depth).
    \item \texttt{min\_data\_in\_leaf}: At least 100--500 for noisy financial data.
    \item \texttt{learning\_rate}: Small (0.01--0.05) with many boosting rounds
          and early stopping.
    \item \texttt{feature\_fraction}: Subsample features per tree (0.5--0.8)
          for regularization.
\end{itemize}

\begin{remark}
    The project already includes LightGBM as a dependency, but no ML-based
    strategy has been implemented yet.  The recommended sequence: build the
    factor evaluation framework (Part~III) first, then use ML to \emph{combine}
    factors, not to \emph{replace} them.
\end{remark}

\begin{labmod}
    \textbf{Build it}: Train a LightGBM model to predict 5-day forward return
    \emph{direction} (binary classification) using the feature matrix from
    Chapter~\ref{ch:features}.  Use \texttt{PurgedWalkForwardCV} from
    Chapter~\ref{ch:ts-validation}.

    \begin{enumerate}
        \item Train on 12 months, predict the next month.  Slide monthly.
        \item Report out-of-sample AUC and accuracy.
        \item Extract feature importance.  Which features matter most?
        \item Compare with a simple equal-weight factor composite (Part~III).
              Does ML add incremental alpha?
    \end{enumerate}
\end{labmod}

% ── Experiment ────────────────────────────────────────────────────────────────
\section{Experiment: LightGBM vs.\ Linear Factor Model}

\begin{experiment}
    \textbf{Goal}: Determine whether ML adds value beyond a linear factor
    composite.

    \textbf{Method}:
    \begin{enumerate}
        \item Baseline: Equal-weight composite of 5+ factors (Part~III result).
              Evaluate with IC/IR.
        \item LightGBM: Same features, walk-forward trained.  Convert
              predicted probabilities to a cross-sectional signal.  Evaluate
              with IC/IR.
        \item Compare IC, IR, and portfolio-level Sharpe.
    \end{enumerate}

    \textbf{Hypothesis}: ML will capture nonlinear interactions and improve IC
    by 10--30\%, but at the cost of lower interpretability and higher turnover.

    \textbf{Reflection space}: \emph{(fill in after running)}
\end{experiment}


% ══════════════════════════════════════════════════════════════════════════════
% APPENDICES
% ══════════════════════════════════════════════════════════════════════════════

\appendix

% ──────────────────────────────────────────────────────────────────────────────
\chapter{Mathematical Reference}
\label{app:math}

This appendix collects the mathematical foundations referenced by the main
chapters.  It is not meant to be read front-to-back --- flip here when a chapter
says ``see Appendix~A.''


\section{Probability and Statistics}
\label{sec:math-stats}

\begin{definition}[Random Variable]
    A measurable function $X: \Omega \to \mathbb{R}$ from a sample space to the
    reals.  In finance, $X$ typically represents a future return or a signal
    value.
\end{definition}

\begin{definition}[Expected Value]
    Discrete: $\mathbb{E}[X] = \sum_x x \cdot p(x)$.
    Continuous: $\mathbb{E}[X] = \int x \, f(x)\, dx$.
\end{definition}

\begin{definition}[Variance and Standard Deviation]
    \[
        \text{Var}(X) = \mathbb{E}[(X - \mathbb{E}[X])^2]
        \qquad
        \sigma(X) = \sqrt{\text{Var}(X)}
    \]
    In finance, $\sigma$ of returns is the standard measure of
    \textbf{volatility}.
\end{definition}

\begin{definition}[Covariance and Correlation]
    \[
        \text{Cov}(X,Y) = \mathbb{E}[(X-\mu_X)(Y-\mu_Y)]
        \qquad
        \rho(X,Y) = \frac{\text{Cov}(X,Y)}{\sigma_X \sigma_Y}
    \]
    Correlation measures \emph{linear} dependence only.
\end{definition}

\begin{definition}[Conditional Expectation]
    $\mathbb{E}[X \mid Y = y]$: the expected value of $X$ given that $Y = y$.
    This is the foundation of alpha research --- ``what is the expected return
    conditional on observing this signal value?''
\end{definition}

\subsection{Key Theorems}

\begin{definition}[Law of Large Numbers]
    As $n \to \infty$: $\bar{X}_n \to \mathbb{E}[X]$.
    \emph{Practical}: Your backtest's average return converges to the true
    expected return --- if the process is stationary.
\end{definition}

\begin{definition}[Central Limit Theorem]
    For i.i.d.\ $X_i$ with mean $\mu$, variance $\sigma^2$:
    \[
        \frac{\bar{X}_n - \mu}{\sigma/\sqrt{n}} \xrightarrow{d} \mathcal{N}(0,1)
    \]
    \emph{Practical}: Justifies $t$-tests on Sharpe ratios and IC values ---
    but financial returns violate i.i.d.\ (autocorrelation, fat tails), so
    apply with caution.
\end{definition}

\subsection{Hypothesis Testing}

\begin{definition}[Hypothesis Test]
    Test $H_0$ (e.g., ``IC $= 0$'') against $H_1$ (``IC $\neq 0$''):
    \begin{enumerate}
        \item Compute a test statistic (e.g., $t = \bar{x} / (s/\sqrt{n})$).
        \item Find the $p$-value: probability of observing this extreme a value
              under $H_0$.
        \item Reject $H_0$ if $p < \alpha$ (typically $\alpha = 0.05$).
    \end{enumerate}
\end{definition}

\subsection{Regression}

\begin{definition}[OLS]
    Minimize $\sum_i (y_i - x_i^T\beta)^2$.  Solution:
    $\hat{\beta} = (X^TX)^{-1}X^Ty$.

    In quant finance: cross-sectional regression of stock returns on factor
    exposures yields \textbf{factor returns} at each date.
\end{definition}

\begin{definition}[Regularization]
    When factors are correlated:
    \begin{itemize}
        \item \textbf{Ridge (L2)}: $+ \lambda\|\beta\|_2^2$.  Shrinks
              coefficients.
        \item \textbf{Lasso (L1)}: $+ \lambda\|\beta\|_1$.  Drives
              irrelevant coefficients to zero (feature selection).
    \end{itemize}
\end{definition}


\section{Time Series}
\label{sec:math-ts}

\begin{definition}[Stationarity]
    A time series $\{X_t\}$ is weakly stationary if
    $\mathbb{E}[X_t] = \mu$ for all~$t$ and
    $\text{Cov}(X_t, X_{t+h})$ depends only on the lag~$h$.
    Prices are non-stationary; returns are approximately stationary.
\end{definition}

\begin{definition}[Unit Root and ADF Test]
    Unit root: $X_t = X_{t-1} + \epsilon_t$ (random walk).  The Augmented
    Dickey--Fuller test tests $H_0$: unit root exists.  Reject $\Rightarrow$
    the series is stationary.
\end{definition}

\begin{definition}[Cointegration]
    Two non-stationary series $X_t$, $Y_t$ are cointegrated if
    $\exists\,\beta$ such that $Y_t - \beta X_t$ is stationary.
    Foundation of pairs trading / statistical arbitrage.
\end{definition}

\begin{definition}[ACF / PACF]
    \begin{itemize}
        \item ACF: $\rho(h) = \text{Corr}(X_t, X_{t+h})$.
        \item PACF: Correlation between $X_t$ and $X_{t+h}$ after removing
              linear effects of intermediate lags.
    \end{itemize}
    Diagnostic tools for identifying AR and MA order.
\end{definition}

\begin{definition}[ARIMA]
    $\text{AR}(p)$: $X_t = \sum_{i=1}^p \phi_i X_{t-i} + \epsilon_t$. \quad
    $\text{MA}(q)$: $X_t = \epsilon_t + \sum_{j=1}^q \theta_j \epsilon_{t-j}$. \quad
    $\text{ARIMA}(p,d,q)$: ARMA on the $d$-th difference.
\end{definition}

\begin{definition}[GARCH(1,1)]
    \[
        \sigma_t^2 = \omega + \alpha\,\epsilon_{t-1}^2 + \beta\,\sigma_{t-1}^2
    \]
    Captures volatility clustering: today's variance depends on yesterday's
    shock and variance.
\end{definition}

\begin{definition}[Realized Volatility]
    $RV_t = \sum_{i=1}^N r_{t,i}^2$ where $r_{t,i}$ are intraday returns
    (e.g., 5-minute).  Uses high-frequency data to estimate daily vol more
    accurately than daily close-to-close.
\end{definition}


\section{Stochastic Calculus}
\label{sec:math-stoch}

These are foundations for derivatives pricing and advanced risk modeling.  Not
immediately needed for equity alpha research, but essential for a complete
quant education.

\begin{definition}[Brownian Motion]
    A continuous-time process $W_t$ with $W_0 = 0$, independent Gaussian
    increments ($W_t - W_s \sim \mathcal{N}(0, t-s)$), and continuous paths.
\end{definition}

\begin{definition}[Geometric Brownian Motion]
    \[
        dS_t = \mu S_t\,dt + \sigma S_t\,dW_t
    \]
    Solution: $S_t = S_0 \exp((\mu - \sigma^2/2)t + \sigma W_t)$.
    The basis of the Black--Scholes model.
\end{definition}

\begin{definition}[It\^{o}'s Lemma]
    For $f(S_t, t)$ where $S_t$ follows an It\^{o} process:
    \[
        df = \frac{\partial f}{\partial t}\,dt
           + \frac{\partial f}{\partial S}\,dS
           + \frac{1}{2}\frac{\partial^2 f}{\partial S^2}(dS)^2
    \]
    The $(dS)^2$ term distinguishes stochastic calculus from ordinary calculus.
\end{definition}

\begin{definition}[Martingale]
    $\{M_t\}$ is a martingale if
    $\mathbb{E}[M_t \mid \mathcal{F}_s] = M_s$ for all $s \leq t$.
    Under the risk-neutral measure, discounted asset prices are martingales.
\end{definition}


% ──────────────────────────────────────────────────────────────────────────────
\chapter{Quant Notes}
\label{app:notes}

This appendix records key insights discovered during development and study.
Each entry captures the date, the question that arose, the finding,
and why it matters.  Entries are ordered chronologically.

% ── Entry 1 ──────────────────────────────────────────────────────────────────
\bigskip\noindent\rule{\textwidth}{0.4pt}

\paragraph{Entry 1 --- Individual Factor IR Is Weak by Design}
\label{exp:weak-ir}
\begin{description}[style=nextline, leftmargin=1.8cm]
    \item[Date] 2026-03-02 \quad (from \texttt{alpha\_iteration.ipynb})
    \item[Question]
        All four factors (BBIBOLL, STR, VPD, Vol~Ratio) showed
        $|\mathrm{IR}|$ in the $0.05$--$0.12$ range.
        Are these signals too weak to be useful?
    \item[Finding]
        No --- this is \emph{normal}.
        In practice, single-factor IRs rarely exceed $0.3$;
        values above $0.5$ are exceptional.
        The practical ranges are roughly:
        \begin{itemize}[nosep]
            \item $|\mathrm{IR}| < 0.05$: noise, not tradeable.
            \item $0.05 \le |\mathrm{IR}| < 0.15$: weak but real ---
                  the typical range for individual alpha factors.
            \item $0.15 \le |\mathrm{IR}| < 0.30$: solid single factor.
            \item $|\mathrm{IR}| \ge 0.50$: suspicious --- likely overfitted
                  or measured in-sample.
        \end{itemize}

    \item[Why It Matters]
        The \textbf{Fundamental Law of Active Management}
        (Grinold \& Kahn) explains how weak signals become useful:
        \[
            \mathrm{IR}_{\text{portfolio}}
            \;\approx\; \mathrm{IR}_{\text{single}} \times \sqrt{N}
        \]
        where $N$ is the number of \emph{independent} bets (signals $\times$
        rebalance periods).
        Combining 10 orthogonal factors each with $\mathrm{IR} = 0.10$ gives
        a portfolio-level $\mathrm{IR} \approx 0.10 \times \sqrt{10} \approx 0.32$,
        well into the ``strong'' range.
        The game is not finding one brilliant factor; it is combining many
        weak-but-orthogonal signals.
\end{description}

% ── Entry 2 ──────────────────────────────────────────────────────────────────
\bigskip\noindent\rule{\textwidth}{0.4pt}

\paragraph{Entry 2 --- Anti-Correlation Requires Aligned IC Signs}
\label{exp:ic-sign-anticorr}
\begin{description}[style=nextline, leftmargin=1.8cm]
    \item[Date] 2026-03-02 \quad (from \texttt{alpha\_iteration.ipynb})
    \item[Question]
        BBIBOLL and STR have strongly anti-correlated daily IC
        ($\rho = -0.80$).
        Anti-correlation should be great for diversification ---
        why did equal-weight combination \emph{dilute} the signal instead?
    \item[Finding]
        Anti-correlation in the IC \emph{series} means the two factors
        ``take turns being right.''
        But whether that helps depends on whether they agree on
        \textbf{direction}:
        \begin{itemize}[nosep]
            \item \textbf{Same IC sign + anti-correlated $\boldsymbol{\to}$ great.}
                  Both factors predict the same direction (e.g.\ both negative IC).
                  On days factor~A is strong, factor~B is weak --- but they still
                  point the same way.  The composite mean IC is preserved while the
                  variance drops, so IR rises.
            \item \textbf{Opposite IC sign + anti-correlated $\boldsymbol{\to}$ cancellation.}
                  BBIBOLL has $\mathrm{IC} = -0.022$ and STR has
                  $\mathrm{IC} = +0.013$.
                  When BBIBOLL fires strongly ($\mathrm{IC} \approx -0.05$),
                  the anti-correlation means STR also fires strongly ---
                  but in the \emph{opposite} direction ($\mathrm{IC} \approx +0.04$).
                  Equal-weight averaging: $(-0.05 + 0.04)/2 \approx -0.005$.
                  The signals cancel, collapsing the mean IC.
        \end{itemize}

    \item[The Math]
        For two equal-weight factors with IC means $\mu_A, \mu_B$,
        IC standard deviations $\sigma_A, \sigma_B$, and IC correlation $\rho$:
        \begin{align*}
            \mu_{\text{combo}}   &= \tfrac{1}{2}(\mu_A + \mu_B) \\
            \sigma^2_{\text{combo}} &= \tfrac{1}{4}
                \bigl(\sigma_A^2 + \sigma_B^2 + 2\rho\,\sigma_A\sigma_B\bigr)
        \end{align*}
        Anti-correlation ($\rho < 0$) always shrinks $\sigma^2_{\text{combo}}$
        (the denominator of IR).
        But if $\mu_A$ and $\mu_B$ have opposite signs,
        $\mu_{\text{combo}}$ also shrinks (the numerator of IR).
        The net effect on $\mathrm{IR} = \mu/\sigma$ depends on which
        shrinks faster.  In the BBIBOLL $\times$ STR case, the numerator
        collapsed more than the denominator --- so IR fell.

    \item[Experimental Evidence]
        \begin{itemize}[nosep]
            \item BBIBOLL + Vol~Ratio ($\rho = 0.02$, same IC sign):
                  composite $|\mathrm{IR}| = 0.136 > 0.122$ (best individual).
                  \textbf{Diversification confirmed.}
            \item BBIBOLL + STR ($\rho = -0.80$, opposite IC sign):
                  composite IR diluted.  Anti-correlation \emph{hurt}.
        \end{itemize}

    \item[Why It Matters]
        The analogy is portfolio construction: two negatively correlated
        assets reduce portfolio volatility only if you are \textbf{long both}.
        If you are long one and short the other, negative correlation means
        they move together --- more risk, not less.
        Similarly, combining factors with opposite IC signs is like going
        long one alpha and short the other; the anti-correlation amplifies
        cancellation rather than providing stability.

        \medskip
        \textbf{Interview one-liner:}
        \emph{``Anti-correlated ICs reduce tracking error, but only improve IR
        if the factors agree on direction.
        Opposite IC signs turn anti-correlation from a hedge into
        a cancellation.''}
\end{description}

% ── Entry 3 ──────────────────────────────────────────────────────────────────
\bigskip\noindent\rule{\textwidth}{0.4pt}

\paragraph{Entry 3 --- Strategy-Based vs.\ Alpha-Pipeline Backtesting}
\label{exp:strategy-vs-pipeline}
\begin{description}[style=nextline, leftmargin=1.8cm]
    \item[Date] 2026-03-02 \quad (architecture comparison)
    \item[Question]
        The codebase offers two backtest paths:
        \begin{itemize}[nosep,topsep=2pt]
            \item \texttt{BBIBOLLStrategy} $\to$
                  \texttt{BacktestEngine}
                  (in \texttt{src/strategy/})
            \item \texttt{pipeline.py} $\to$
                  \texttt{WeightBacktester}
                  (in \texttt{src/portfolio/})
        \end{itemize}
        What is the fundamental difference, and when should
        each be used?

    \item[Finding]
        They embody two distinct mental models of systematic investing.

        \medskip
        \textbf{Strategy approach --- ``I Trade.''}
        \begin{enumerate}[nosep]
            \item Calculate indicators (BBI + Bollinger Bands) per ticker.
            \item Generate signals by iterating \emph{row by row},
                  tracking stateful holding logic via
                  \texttt{is\_holding} and \texttt{buy\_price}.
                  Emit buy$(+1)$/sell$(-1)$ on threshold triggers.
            \item Match buy$\to$sell pairs via \texttt{join\_asof};
                  compute per-trade P\&L.
            \item Expand trades into daily holdings; equal-weight
                  average $\to$ equity curve.
        \end{enumerate}
        Key properties: \emph{stateful} (sell decision depends on entry
        price), \emph{discrete events} (in or out), \emph{single-asset
        logic} (each ticker processed independently),
        \emph{trade-centric output} (``bought AAPL at \$150,
        sold at \$165, return $= +10\%$'').

        \medskip
        \textbf{Pipeline approach --- ``I Allocate.''}
        \begin{enumerate}[nosep]
            \item Compute factors (bbiboll deviation, volatility ratio,
                  momentum) for \emph{all} stocks simultaneously.
            \item Rank cross-sectionally on each rebalance date;
                  assign portfolio weights (top-ranked $\to$ equal weight,
                  rest $\to$ zero).
            \item Portfolio return:
                  $r_t = \sum_i w_{i,t}\, r_{i,t}$
                  --- no explicit trades, just weights $\times$ returns.
        \end{enumerate}
        Key properties: \emph{stateless} (today's weight depends only on
        today's factor values), \emph{continuous allocation} (how much,
        not whether), \emph{cross-sectional} (all stocks compete for
        capital), \emph{portfolio-centric output} (daily portfolio
        return, Sharpe ratio).

    \item[Summary Table]~\\[-\baselineskip]
        \begin{tabular}{@{}lll@{}}
        \toprule
        \textbf{Dimension} & \textbf{Strategy} & \textbf{Pipeline} \\
        \midrule
        Unit of analysis & Individual trade   & Portfolio weight \\
        State            & Entry price, holding flag & None (recomputed fresh) \\
        Decision         & Buy/sell \emph{this} stock? & Allocate across \emph{all} stocks? \\
        Output           & Trade list + P\&L   & Weight matrix + return series \\
        Rebalance        & Signal-triggered    & Calendar-based (weekly/monthly) \\
        Typical role      & Trader / execution  & Researcher / PM \\
        \bottomrule
        \end{tabular}

    \item[Why It Matters]
        Most quant interview questions and institutional workflows assume
        the \textbf{pipeline worldview}: ``describe an alpha factor'',
        ``construct a long-short portfolio'', ``what is your Sharpe?''
        The strategy approach is valuable for execution-level or
        high-frequency work, but the pipeline is closer to how
        firms like Citadel, Two~Sigma, and AQR structure research.

        \medskip
        The bridge between the two is \texttt{WeightBacktester}:
        it accepts pipeline output~(weights) and produces the equity
        curve $+$ Sharpe that connect the ``I allocate'' mindset to
        concrete performance metrics.  In the codebase,
        \texttt{backtester.py --mode strategy} runs the trade-level path
        and \texttt{--mode pipeline} runs the portfolio-level path,
        allowing direct comparison on the same data.

        \medskip
        \textbf{Learning path:} start with the strategy approach (it is
        intuitive --- you can inspect individual trades), then graduate
        to the pipeline (it forces you to think in portfolios:
        factors $\to$ ranking $\to$ allocation $\to$ rebalancing).
        The key insight from the \textbf{Fundamental Law of Active
        Management} (Entry~\ref{exp:weak-ir}) is that the pipeline
        naturally maximises the number of independent bets~$N$, which
        is how weak individual signals become a strong portfolio.
\end{description}

% ── Entry 4 ──────────────────────────────────────────────────────────────────
\bigskip\noindent\rule{\textwidth}{0.4pt}

\paragraph{Entry 4 --- Scalar Kelly vs.\ Multivariate Kelly: Layer Matters}
\label{exp:kelly-lesson}
\begin{description}[style=nextline, leftmargin=1.8cm]
    \item[Date] 2026-03-02 \quad (from \texttt{pipeline\_demo.ipynb} ablation study;
          revised 2026-03-05)

    \item[Question]
        We implemented ``Half-Kelly'' position sizing:
        $w_i \propto \mu_i / \sigma_i^2$, normalized to unit gross leverage.
        Every configuration using it produced \emph{worse} Sharpe than
        equal-weight.  Three successive bug fixes failed to improve it.
        What is fundamentally wrong?

    \item[Background --- The Kelly Formula]
        The Kelly criterion solves: \emph{``What fraction of my bankroll
        should I wager on a bet with known edge $\mu$ and variance
        $\sigma^2$?''}
        \[
            f^* = \frac{\mu}{\sigma^2}
        \]
        This is optimal for a \textbf{single repeated bet} with
        independent outcomes, maximizing the long-run geometric growth
        rate $G = \mathbb{E}[\log(1 + f \cdot r)]$.

        Half-Kelly ($f^*/2$) is a common practitioner adjustment that
        sacrifices roughly $25\%$ of growth for substantially lower
        drawdown.  The key property: \emph{Kelly tells you how much
        total capital to risk}, not how to split capital across
        multiple simultaneous bets.

    \item[The Three Bugs We Hit (and Fixed)]
        \begin{enumerate}[nosep]
            \item \textbf{Normalization destroyed the leverage signal.}
                  We normalized $\sum|w_i| = 1$, which discards Kelly's
                  answer to ``how much to bet.''  After normalization,
                  only the \emph{relative} sizing remains --- but that
                  was never Kelly's purpose.
                  \emph{Fix:} let leverage float (partly undone by
                  needing comparable backtest results).
            \item \textbf{Using $\mu_i$ from returns made Kelly ignore
                  the factor.}
                  Historical mean return over 60 days is dominated by
                  noise (SNR $\approx 0.02$); the factor signal was
                  unused.  All \texttt{AlphaConfig} variants produced
                  identical Sharpe $= 0.343$ regardless of factor
                  choice.
                  \emph{Fix:} two-stage design --- factor selects
                  stocks, Kelly sizes by $|\mu_i| / \sigma_i^2$.
            \item \textbf{Historical $\mu$ is stealth momentum.}
                  Even after the two-stage fix, $\mu_i$ from a 60-day
                  window injects a momentum bias.  For a mean-reversion
                  signal like BBIBOLL, this conflicts with the factor
                  direction.
                  \emph{Fix:} replace $\mu_i$ with $z_i$, the
                  cross-sectional z-score of the factor signal itself.
        \end{enumerate}

    \item[The Ablation Study]
        After all three fixes, we ran a $2 \times 2$ experiment
        isolating sizing and rebalancing frequency:
        \begin{center}
        \begin{tabular}{@{}llcc@{}}
        \toprule
        & & \multicolumn{2}{c}{\textbf{Rebalance}} \\
        \cmidrule(l){3-4}
        & & Daily & Weekly \\
        \midrule
        \multirow{2}{*}{\textbf{Sizing}}
        & Equal-Weight      & $+0.814$ & $+0.417$ \\
        & Signal-Weighted   & $+0.137$ & $-0.393$ \\
        \bottomrule
        \end{tabular}
        \end{center}

        \medskip\noindent
        Two independent sources of Sharpe degradation:
        \begin{itemize}[nosep]
            \item \textbf{Weekly rebalancing} $\to$ $-0.40$ Sharpe
                  (mean-reversion alpha decays intraweek).
            \item \textbf{Signal-weighted sizing} $\to$ $-0.68$ Sharpe
                  (concentration penalty).
        \end{itemize}

    \item[Root Cause --- Why ``Kelly Sizing'' Hurts]
        The $z_i / \sigma_i^2$ formula concentrates weight into
        low-volatility, high-conviction names.  But for BBIBOLL
        (a mean-reversion signal), the strongest alpha comes from
        \emph{high-volatility} stocks that have deviated far from
        their bands.  Signal-weighted sizing systematically
        \textbf{under-weights the best opportunities}.

        \medskip
        Mathematically: $w_i \propto z_i / \sigma_i^2$ penalizes
        large $\sigma_i$ quadratically.  In cross-section, the stocks
        with the largest $|z_i|$ (strongest signal) also tend to have
        the largest $\sigma_i$ (that is what drove the deviation).
        The penalty overwhelms the conviction signal.

    \item[The Conceptual Error]
        Kelly solves a fundamentally different problem than what
        we asked it to do:
        \begin{center}
        \begin{tabular}{@{}lp{5.5cm}p{5.5cm}@{}}
        \toprule
        & \textbf{Kelly's Actual Problem} & \textbf{What We Used It For} \\
        \midrule
        Question & How much total capital to deploy?
                 & How to split capital across $N$ stocks? \\
        Input    & Edge and variance of \emph{one} bet
                 & Cross-sectional factor signal \\
        Output   & Portfolio-level leverage $f^*$
                 & Relative stock weights $w_i$ \\
        Regime   & Independent repeated bets
                 & Correlated concurrent positions \\
        \bottomrule
        \end{tabular}
        \end{center}

        \medskip
        In a multi-asset portfolio, the correct Kelly generalization
        is the \textbf{multivariate Kelly}:
        \[
            \mathbf{f}^* = \Sigma^{-1}\,\boldsymbol{\mu}
        \]
        where $\Sigma$ is the covariance matrix and $\boldsymbol{\mu}$
        is the expected return vector.  This accounts for
        cross-correlations and produces a single leverage-scaled
        portfolio.  Our per-stock $\mu_i / \sigma_i^2$ ignores
        correlations entirely.

    \item[Resolution]
        \begin{enumerate}[nosep]
            \item \textbf{Renamed} \texttt{size\_half\_kelly} $\to$
                  \texttt{size\_signal\_weighted} with docstring
                  explicitly stating ``Not true Kelly.''
            \item \textbf{Formula kept} as
                  $w_i \propto \mathrm{direction}_i \times |z_i| / \sigma_i^2$
                  --- it is a valid conviction~$\times$~inverse-variance
                  sizing method, just not Kelly.
            \item \textbf{Multivariate Kelly belongs at Layer~2
                  (portfolio construction).}
                  The correct multi-asset generalization
                  $\mathbf{f}^* = \Sigma^{-1}\boldsymbol{\mu}$
                  \emph{produces weights}, not just a scalar leverage.
                  Half-Kelly ($\mathbf{f}^*/2$) is the practical
                  version.  This is a legitimate portfolio
                  construction method --- but it lives at the
                  \emph{strategy backtest} layer, not the factor
                  evaluation layer.  At the factor evaluation layer
                  (alpha research), equal-weight long-short is the
                  standard.
        \end{enumerate}

        \medskip\noindent
        \textbf{Where each Kelly variant belongs:}
        \begin{center}
        \small
        \resizebox{\linewidth}{!}{%
        \begin{tabular}{@{}lll@{}}
        \toprule
        \textbf{Variant} & \textbf{Formula} & \textbf{Layer} \\
        \midrule
        Scalar ``Kelly'' ($\mu_i/\sigma_i^2$ per stock)
            & Ignores correlations
            & Misapplication --- not valid at any layer \\
        Signal-weighted ($z_i/\sigma_i^2$, normalized)
            & Conviction $\times$ inverse-var
            & Layer~2: position sizing \\
        Multivariate Kelly ($\Sigma^{-1}\mu$)
            & Full covariance-aware
            & Layer~2: portfolio construction \\
        \bottomrule
        \end{tabular}}%
        \end{center}

    \item[Why It Matters]
        \begin{enumerate}[nosep]
            \item \textbf{Naming matters.}  Calling something ``Kelly''
                  imports assumptions (independent bets, known $\mu$)
                  that do not hold.  Honest names prevent honest mistakes.
            \item \textbf{Position sizing is not ``just implementation.''}
                  With the \emph{same} signal, Sharpe ranged from
                  $-0.39$ to $+0.81$ --- a 1.2 swing.  Sizing is
                  a core part of the strategy, not a post-hoc detail.
            \item \textbf{Kelly requires reliable $\mu$ and $\Sigma$
                  estimates.}
                  With $|\mathrm{IC}| \approx 0.02$, the expected
                  return signal is noise.  Scalar Kelly amplifies
                  noise; multivariate Kelly adds $\Sigma^{-1}$
                  estimation error (covariance inversion is
                  notoriously unstable for large $N$).  Both become
                  useful only with $\mathrm{IC} > 0.05$ and
                  well-conditioned covariance estimates.
        \end{enumerate}

        \medskip
        \textbf{Interview one-liner:}
        \emph{``Scalar Kelly ignores correlations --- it's a
        misapplication.  Multivariate Kelly
        ($\mathbf{f}^*=\Sigma^{-1}\mu$) is a legitimate portfolio
        construction method, but it belongs at the strategy layer,
        not the alpha research layer.  At the research layer,
        equal-weight long-short is the standard.''}
\end{description}

\bigskip\noindent\rule{\textwidth}{0.4pt}


% ── Entry 5 ──────────────────────────────────────────────────────────────────
\bigskip\noindent\rule{\textwidth}{0.4pt}

\paragraph{Entry 5 --- Position Sizing Is Portfolio Construction, Not Just Weighting}
\label{exp:position-sizing-landscape}
\begin{description}[style=nextline, leftmargin=1.8cm]
    \item[Date] 2026-03-04 \quad (revised 2026-03-05)

    \item[Context]
        After renaming Half-Kelly $\to$ Signal-Weighted (Entry~4),
        we stepped back to examine the position sizing module as a
        whole.  All four methods share the same two-stage architecture:

        \begin{enumerate}[nosep]
            \item \textbf{Selection} --- rank stocks by signal, pick
                  top-$N$ long and bottom-$N$ short (binary in/out).
            \item \textbf{Sizing} --- within each leg, assign weights
                  (equal, inverse-vol, vol-target, or signal-weighted).
        \end{enumerate}

        This is the standard \emph{rank-and-select} approach from
        academic factor research (Fama--French quintile portfolios).

    \item[The Long-Short Confusion (and Correction)]
        We initially believed that setting \texttt{n\_short=10} was a
        \emph{design bug} --- forcing the portfolio to short stocks
        that the signal merely ranked lower, not genuinely bearish.
        We added \texttt{portfolio\_mode} to \texttt{AlphaConfig} and
        switched the default to \texttt{"long\_only"}.

        \medskip
        \textbf{What we missed:}  Long-short evaluation is not a bug
        --- it is the \emph{standard methodology} for factor research.
        The purpose is to construct a \textbf{market-neutral} portfolio
        that isolates the pure cross-sectional factor return (alpha)
        from the market return (beta).  If you only go long the top-20
        stocks and the market rallies 30\%, you cannot tell whether the
        alpha is real or the market carried you.

        \medskip
        \textbf{Three layers, three correct answers:}

        \begin{center}
        \small
        \resizebox{\linewidth}{!}{
        \begin{tabular}{@{}llp{7cm}@{}}
        \toprule
        \textbf{Layer}  & \textbf{Mode} & \textbf{Why} \\
        \midrule
        1. Factor evaluation
            & Always L/S
            & Isolate alpha from beta.  Quintile spread, IC, factor
              returns.  Standard academic method. \\
        2. Strategy backtest
            & Depends on mandate
            & Hedge fund $\to$ long-short.  Long-only fund or retail
              $\to$ set \texttt{n\_short=0}.  This is where
              \texttt{portfolio\_mode} lives. \\
        3. Production execution
            & Separate system
            & Execution algos, slippage, margin, order routing.
              Not in scope yet. \\
        \bottomrule
        \end{tabular}}
        \end{center}

        \medskip
        \textbf{The \texttt{portfolio\_mode} addition is still
        correct} --- it prevents accidentally forcing shorts when
        backtesting a long-only deployment strategy (layer~2).  The
        mistake was \emph{conflating} layer~1 (evaluation) with
        layer~2 (backtesting).  At layer~1, long-short is not
        ``forced shorting'' --- it is intentional isolation of the
        factor spread.

    \item[Lesson]
        \begin{quote}
        \emph{When your factor evaluation shows Long~Q1, Short~Q5,
        and someone says ``why are we shorting good stocks?''\ ---
        the answer is: we are not betting against them.  We are
        measuring whether the \textbf{ranking} has predictive power,
        independent of market direction.  The short leg is a
        research tool, not a trading decision.}
        \end{quote}

    \item[Why Fixed $N$?]
        Fixed $N$ has real advantages:
        \begin{enumerate}[nosep]
            \item \textbf{Stable risk profile} --- constant gross
                  exposure, predictable diversification and capacity.
            \item \textbf{No threshold to tune} --- avoids the
                  hyperparameter ``what signal level is good enough?''
            \item \textbf{Apples-to-apples} --- isolates sizing effect
                  from selection effect when comparing methods.
            \item \textbf{Academic convention} --- standard in factor
                  evaluation (quintile L/S, decile spreads).
        \end{enumerate}

        \medskip
        But it also has fundamental limitations:
        \begin{itemize}[nosep]
            \item Stock ranked \#11 gets zero weight regardless of
                  signal strength (information discarded at selection).
            \item Forced positions when fewer than $N$ stocks have
                  genuinely positive signal.
            \item Signal magnitude is lost at the selection stage
                  (Signal-Weighted partially recovers it at sizing).
        \end{itemize}

    \item[The Portfolio Construction Spectrum]
        The rank-and-select approach is one end of a spectrum.  The
        other end is \textbf{optimization-based} construction, where
        selection and sizing happen simultaneously:

        \medskip
        \begin{center}
        \small
        \begin{tabular}{@{}lll@{}}
        \toprule
        \textbf{Approach} & \textbf{Idea} & \textbf{Tradeoff} \\
        \midrule
        Equal-Weight   & $1/N$ each position      & Robust; ignores signal \& risk \\
        Inverse-Vol    & $w_i \propto 1/\sigma_i$  & Equal risk; needs $\sigma$ estimates \\
        Vol-Target     & Scale to $\sigma_{\text{target}}$ & Leverage varies; chases past vol \\
        Signal-Weighted & $w_i \propto |z_i|/\sigma_i^2$ & Uses conviction; ignores correlations \\
        \midrule
        \textbf{MVO}   & $w^* = \arg\max (w^\top\mu - \frac{\lambda}{2} w^\top\Sigma w)$ & Optimal but fragile \\
        Risk Parity    & Equal risk contribution   & Robust; ignores $\mu$ \\
        Black-Litterman & Bayesian prior + views   & Principled; complex \\
        \bottomrule
        \end{tabular}
        \end{center}

    \item[Mean-Variance Optimization (MVO)]
        The closed-form solution is:
        \[
            w^* = \frac{1}{\lambda}\,\Sigma^{-1}\,\mu
        \]
        where $\mu$ is the expected return vector (from the alpha
        signal), $\Sigma$ is the covariance matrix, and $\lambda$
        is risk aversion.

        \medskip
        \textbf{Key insight}: MVO solves selection and sizing
        \emph{simultaneously}.  Every stock gets a weight proportional
        to its signal strength, adjusted for correlation with the rest.
        No fixed $N$, no arbitrary threshold.

        \medskip
        \textbf{Practical problems with MVO}:
        \begin{enumerate}[nosep]
            \item \textbf{Covariance estimation} --- 50 stocks
                  $\Rightarrow$ 1,225 parameters in $\Sigma$.
                  Mitigations: Ledoit--Wolf shrinkage, factor models.
            \item \textbf{Extreme weights} --- $\Sigma^{-1}$ amplifies
                  estimation error.  Must add constraints
                  ($|w_i| \leq w_{\max}$, leverage cap).
            \item \textbf{Sensitivity to $\mu$} --- small changes in
                  expected return $\Rightarrow$ completely different
                  portfolio.  Black-Litterman addresses this.
            \item \textbf{Turnover explosion} --- optimal portfolio
                  changes daily.  Add turnover penalty
                  $\gamma \|w_t - w_{t-1}\|_1$ to the objective.
        \end{enumerate}

        \medskip
        The classic quant aphorism: \emph{``MVO is the world's best
        error maximizer''} --- it takes the noisiest input ($\mu$) and
        amplifies it through the most unstable operation ($\Sigma^{-1}$).

    \item[What This Means for Our Project]
        Our Phase~3 cost analysis proved the BBIBOLL + Vol Ratio alpha
        has a breakeven cost of $\sim$2\,bps.  MVO would not fix this
        --- it would just allocate noise more ``efficiently.''

        \medskip
        \textbf{Priority order}:
        \begin{enumerate}[nosep]
            \item \textbf{Alpha improvement} (more factors, better
                  combination) --- fixes the real bottleneck.
            \item \textbf{MVO implementation} --- becomes worthwhile
                  once the composite signal has breakeven $>$10\,bps.
            \item \textbf{Turnover-aware optimization} --- internalize
                  costs in the objective rather than analyzing post-hoc.
        \end{enumerate}

    \item[Interview Takeaway]
        \emph{``Position sizing is portfolio construction, which sits
        between alpha generation and execution.  Rank-and-select is
        robust but throws away information.  MVO is theoretically
        optimal but fragile --- it amplifies estimation error in
        $\mu$ through $\Sigma^{-1}$.  The right approach depends on
        signal quality: with IC $\approx$ 0.02, robust methods
        dominate; with IC $>$ 0.05, optimization starts to pay off.
        Either way, no portfolio construction method can rescue a
        weak alpha.''}
\end{description}

\bigskip\noindent\rule{\textwidth}{0.4pt}


% ──────────────────────────────────────────────────────────────────────────────
\chapter{Glossary}
\label{app:glossary}

\begin{table}[ht]
\centering
\small
\begin{tabular}{@{}ll@{}}
\toprule
\textbf{Term} & \textbf{Definition} \\
\midrule
ADF       & Augmented Dickey--Fuller test for stationarity \\
BBI       & Bull and Bear Index (average of multiple SMAs) \\
CAGR      & Compound Annual Growth Rate \\
CVaR      & Conditional Value at Risk (Expected Shortfall) \\
FIGI      & Financial Instrument Global Identifier \\
GARCH     & Generalized Autoregressive Conditional Heteroskedasticity \\
GICS      & Global Industry Classification Standard \\
IC        & Information Coefficient \\
IR        & Information Ratio \\
MDD       & Maximum Drawdown \\
MVO       & Mean-Variance Optimization (Markowitz) \\
OBV       & On-Balance Volume \\
OHLCV     & Open, High, Low, Close, Volume \\
OLS       & Ordinary Least Squares \\
RV        & Realized Volatility \\
VaR       & Value at Risk \\
VWAP      & Volume-Weighted Average Price \\
WFO       & Walk-Forward Optimization \\
\bottomrule
\end{tabular}
\caption{Common abbreviations used in this document}
\end{table}


% ──────────────────────────────────────────────────────────────────────────────
\chapter{Learning Priority and Reading List}
\label{app:priority}

\section{Recommended Order}

\begin{enumerate}
    \item \textbf{Data Pipeline} (Part~I): Get clean data flowing first.
          Nothing else works without it.
    \item \textbf{Backtest \& Bias} (Part~II): Build a strategy, measure it,
          understand why the measurement might be wrong.
    \item \textbf{Factor IC/IR Analysis} (Part~III, Ch.~\ref{ch:factor-eval}):
          The core language of quant interviews.
    \item \textbf{Factor Construction} (Part~III, Ch.~\ref{ch:factor-construction}):
          Normalization, neutralization --- the preprocessing that makes or
          breaks a factor.
    \item \textbf{Execution \& Costs} (Part~IV, Ch.~\ref{ch:execution}):
          Makes your backtest realistic.
    \item \textbf{Machine Learning} (Part~V): Only after the linear factor
          framework is solid.
    \item \textbf{Stochastic Calculus} (App.~\ref{sec:math-stoch}): For
          derivatives and advanced risk; not urgent for equity alpha.
\end{enumerate}

\section{Suggested Reading}

\begin{itemize}
    \item Qian, Hua \& Sorensen, \textit{Quantitative Equity Portfolio
          Management} --- the standard reference for factor investing.
    \item Grinold \& Kahn, \textit{Active Portfolio Management} --- the
          ``Fundamental Law of Active Management'' and IC/IR framework.
    \item L\'{o}pez de Prado, \textit{Advances in Financial Machine Learning}
          --- purged cross-validation, meta-labeling, feature importance.
    \item Chan, \textit{Quantitative Trading} --- practical guide for
          independent quant traders.
    \item Shreve, \textit{Stochastic Calculus for Finance I \& II} --- the
          mathematical foundations if you want derivatives pricing.
\end{itemize}


\end{document}
